% !TeX root = ../../Book1.tex
% 纲要标题：## 第19章*　紧 Lie 群的极大环面与根几何
% 纲要位置：几何学计划纲要.md，第 842 行。
% 正文状态：初稿；尚未正式审核。

% 纲要内容开始
% 
% > 本章目标与证明范围：保留紧连通 Lie 群极大环面存在、元素共轭入环面及环面共轭性的完整证明。证明所需的紧 Riemann 引理在本章自备；完整根系分类不作为此前流形理论的前提。
% 
% 纲要内容结束

\OptionalChapter{紧 Lie 群的极大环面与根几何}{g1:ch:19}
\BookChapterNumber{g1:ch:19}

\begin{chapterabstract}
	紧连通 Lie 群中的单个元素虽然未必与其他元素交换，却总能放进一个连通交换的闭子群。这样的极大环面彼此共轭，使群内的共轭问题归结为环面上的有限对称。本章先由归一化积分构造双不变度量，补足紧群上最短路径与指数满射所需的论证，再证明极大环面的存在、元素入环面及环面共轭定理。伴随作用的权分解给出正则元素、根方向和反射，整数格则补充无穷小结构不能记录的整体信息。最后在酉群、旋转群和紧辛群中计算环面、根、Weyl 群、中心与基本群，并用低维覆盖比较相同 Lie 代数的不同群。
\end{chapterabstract}

\begin{learninggoals}
	\begin{itemize}
		\item 从紧群不变积分构造伴随不变内积，说明双不变度量中的一参数子群为何是测地线，并证明指数映射满射。
		\item 由一参数子群的闭包得到环面，证明紧连通交换群的格商结构，并把任一元素放进极大环面。
		\item 用环面作用的无穷小分解辨认正则元素，证明极大环面共轭，区分群元素共轭与 Lie 代数元素共轭。
		\item 解释正规化子、中心化子、根空间与根反射的作用，辨认根格、特征标及指数核分别保存的信息。
		\item 计算经典紧群中的具体环面与 Weyl 群，用中心、基本群和二重覆盖检验局部 Lie 代数与整体群的区别。
	\end{itemize}
\end{learninggoals}

一个酉矩阵可以在某组酉基下写成对角矩阵，但两个酉矩阵通常不能在同一组基下同时对角化。这里已经隐含两个层次的问题：单个元素可以移到怎样的交换子群中，以及不同选择的交换子群之间有什么关系。三维旋转给出另一个图像：每个旋转都有一条轴，绕固定轴的全部旋转组成圆周群；改变轴则通过整体旋转实现。极大环面定理把这两种经验推广到一般紧连通 Lie 群。

\ProseParagraphBreak

从矩阵模型猜测结论并不困难，困难在于一般 Lie 群没有事先给定的特征向量。我们将先用紧性取得伴随不变内积，再沿一条最短路径把任一群元素写成一个指数。该一参数子群的闭包是紧连通交换群，因此成为环面。环面的共轭性则来自另一个极值问题：把伴随轨道上的一个线性函数取到最大，极值条件会迫使两个 Lie 代数元素交换。读者可以把这两个极值步骤分开记忆：前者寻找通往给定群元素的路径，后者改变无穷小方向的基准位置。

\ProseParagraphBreak

本章的 Lie 群均为有限维实 Lie 群，默认的流形条件仍为无边界、Hausdorff、第二可数；涉及主结论时明确列出紧性与连通性。主要准备是\chapref{g1:ch:16}的指数、闭子群与同态积分，\chapref{g1:ch:17}的齐性空间与紧群平均，以及\chapref{g1:ch:18}的主丛水平提升。最短路径所需的度量论证在本章建立。经典群基本群的计算还使用第五章的覆盖与 van Kampen 定理；相应论证会恢复实际需要的版本。

\ProseParagraphBreak

环面记为 \(T\)，其 Lie 代数记为 \(\mathfrak t\)，与切丛的正体记号 \(\mathrm{T}M\) 区分。根采用实线性泛函 \(\alpha\in\mathfrak t^*\)：根向量 \(Z\) 满足 \([H,Z]=\mathrm{i}\alpha(H)Z\)，所以 \(\operatorname{Ad}_{\exp H}Z=e^{\mathrm{i}\alpha(H)}Z\)。由此，指数核 \(\Lambda=\ker(\exp|_{\mathfrak t})\) 与根之间的整数条件是 \(\alpha(\Lambda)\subseteq2\pi\mathbb Z\)。保留这项归一化，在比较圆周角度、根和特征标时就不会遗漏 \(2\pi\)。

\ProseParagraphBreak

延伸阅读可查 Hall 的《Lie Groups, Lie Algebras, and Representations》中题为“Compact Lie Groups and Maximal Tori”的一章\cite{Hall2015}。本章以极大环面定理及具体根模型为范围；一般根系的分类及最高权表示理论需要进一步的代数工具。

% !TeX root = ../../Book1.tex
% 纲要标题：### 19.1 不变积分与双不变度量
% 纲要位置：几何学计划纲要.md，第 846 行。
% 正文状态：初稿；尚未正式审核。

% 纲要内容开始
% 
% * 紧 Lie 群体积的归一化
% * 左右不变性与平均
% * 伴随不变内积和双不变度量
% * 紧连通 Lie 群的最短测地线与指数满射
% 
% 纲要内容结束

\section{不变积分与双不变度量}
\label{g1:sec:1901}

在一般连通 Lie 群中，一个元素可以写成有限个指数的乘积，却未必是单个指数。紧群上的长度提供了进一步的控制：若能找到从单位元到给定元素的最短曲线，并证明这条曲线来自单参数子群，就能把有限乘积改为一个指数。为此，我们需要一个与左右平移同时相容的度量。本节先把第17章的紧群积分用于伴随表示，再从指数微分公式出发建立所需的最短曲线论证。

\subsection{紧 Lie 群体积的归一化}
\label{g1:subsec:1901:01}
设 \(G\) 为有限维实紧 Lie 群，暂不要求连通，\(\mathfrak g=\mathrm{T}_eG\)。\cref{g1:thm:compact-group-integration}给出了 \(G\) 上唯一的归一化左不变光滑正密度 \(\mu_G\)，而且它同时右不变、取逆不变。这里回顾归一化中的两个要点。先把 \(\mathfrak g\) 上任一正密度 \(\nu_e\) 用左平移延拓，得到
\[
	\nu_g(v_1,\ldots,v_n)
	=\nu_e\bigl(\mathrm{d}(L_{g^{-1}})_gv_1,\ldots,
	\mathrm{d}(L_{g^{-1}})_gv_n\bigr),\qquad n=\dim G.
\]
紧性保证 \(0<\int_G\nu<\infty\)，因而可以令 \(\mu_G=\nu/\int_G\nu\)。单位元处的密度空间是一维，所有左不变正密度只相差一个正常数；除以总量后便与最初选择无关。

\ProseParagraphBreak

所谓体积归一化，归一化的是整个群的总量，而非每个坐标片的 Euclid 体积。它也不依赖定向：正密度在换坐标时使用 Jacobian 行列式的绝对值。即使 \(G\) 不连通，上式仍沿左平移规定每个分支上的密度。紧 Lie 群只有有限个连通分支，因为分支是开集，而所有分支构成一个开覆盖；左平移互换这些分支，所以每个分支的归一化体积相同。

\begin{example}\label{g1:exm:normalized-volume-torus}
	在 \(T^r=(S^1)^r\) 上，用角坐标 \(z_j=e^{\mathrm{i}t_j}\) 表示积分，则
	\[
		\int_{T^r}f\,\mathrm{d}\mu_{T^r}
		=\frac1{(2\pi)^r}\int_{[0,2\pi]^r}
		f(e^{\mathrm{i}t_1},\ldots,e^{\mathrm{i}t_r})\,\mathrm{d}t_1\cdots\mathrm{d}t_r.
	\]
	任一常系数正定内积给出的体积密度只是角坐标密度的正常数倍，故归一化后都得到这一积分。相反，加法群 \(\mathbb R^r\) 在 \(r>0\) 时的左不变正密度总量无限，不能作此归一化。归一化积分的存在与任意长度单位的选择是不同问题。
\end{example}

\subsection{左右不变性与平均}
\label{g1:subsec:1901:02}
右不变性是平均伴随表示时实际使用的条件。回顾\cref{g1:thm:compact-group-integration}的证明：左右平移交换，使 \(R_a^*\mu_G\) 仍左不变，故为 \(\Delta(a)\mu_G\)，其中 \(\Delta:G\to\mathbb R_{>0}\) 是连续同态。紧像 \(\log\Delta(G)\) 是 \(\mathbb R\) 的紧加法子群；若含非零数，就含其所有整数倍而无界。因此 \(\Delta\equiv1\)。这一步说明为何我们可以在同一个积分中同时使用左换元和右换元。

\begin{proposition}[表示的平均内积]\label{g1:prop:compact-representation-inner-product}
	设 \(G\) 为有限维实紧 Lie 群，\(\rho:G\to\operatorname{GL}(V)\) 为有限维实光滑表示，\(b\) 为 \(V\) 上的正定内积，\(\mu_G\) 为 \(G\) 上满足 \(\mu_G(G)=1\) 的左右不变 Haar 测度。则
	\[
		\langle u,v\rangle
		=\int_G b(\rho(g)u,\rho(g)v)\,\mathrm{d}\mu_G(g)
	\]
	是 \(\rho(G)\) 不变的正定内积。若 \(V\) 为复向量空间、\(\rho\) 为复线性表示，则平均任一正定 Hermitian 内积也得到不变正定 Hermitian 内积。
\end{proposition}
\begin{proof}
	在固定基下逐分量积分，所得形式仍保留被积式的双线性和对称性。对非零 \(u\)，每个 \(\rho(g)\) 可逆，所以 \(b(\rho(g)u,\rho(g)u)>0\)。\cref{g1:thm:compact-group-integration}的正性给出 \(\langle u,u\rangle>0\)。对 \(a\in G\)，右换元 \(g\mapsto ga\) 给出
	\[
		\langle\rho(a)u,\rho(a)v\rangle
		=\int_G b(\rho(ga)u,\rho(ga)v)\,\mathrm{d}\mu_G(g)
		=\langle u,v\rangle.
	\]
	复情形逐实部、虚部积分，半双线性、Hermitian 对称性和正性由相同论证得到。
\end{proof}

特别地，取 \(V=\mathfrak g\)、\(\rho=\operatorname{Ad}\)，就得到对所有 \(a\in G\) 满足
\[
	\langle\operatorname{Ad}_aX,\operatorname{Ad}_aY\rangle
	=\langle X,Y\rangle
\]
的\term[bansui-bubian-neiji]{伴随不变内积}[Ad-invariant inner product]。若只在 \(G^0\) 上平均，得到的内积未必被其余分支保持；对整个紧群积分消除了这个问题。

\subsection{伴随不变内积和双不变度量}
\label{g1:subsec:1901:03}
左平移能把任意 \(\mathfrak g\) 上的内积延拓成左不变度量，但右平移会在单位元数据上引入共轭。伴随不变性恰好使这两种要求相容。

\begin{proposition}[双不变度量]\label{g1:prop:bi-invariant-metric}
	设 \(G\) 为有限维实 Lie 群，\(\mathfrak g=\mathrm{T}_eG\)，\(\langle\ ,\ \rangle\) 为 \(\mathfrak g\) 上的正定内积。定义
	\[
		h_g(v,w)=\left\langle\mathrm{d}(L_{g^{-1}})_gv,
		\mathrm{d}(L_{g^{-1}})_gw\right\rangle.
	\]
	则 \(h\) 是光滑左不变 Riemann 度量；它同时右不变，当且仅当该内积 \(\operatorname{Ad}(G)\) 不变。每个紧 Lie 群因而具有\term[shuangbubian-duliang]{双不变度量}[bi-invariant metric]。若 \(G\) 连通，则伴随不变性又等价于
	\begin{equation}\label{g1:eq:ad-skew-inner-product}
		\langle[Z,X],Y\rangle+\langle X,[Z,Y]\rangle=0
		\qquad(X,Y,Z\in\mathfrak g).
	\end{equation}
\end{proposition}
\begin{proof}
	平移微分光滑且可逆，故上式光滑、正定。\(L_{(ag)^{-1}}\circ L_a=L_{g^{-1}}\) 说明左不变性。另一方面，在 \(g\) 处求微分有
	\[
		\mathrm{d}(L_{(ga)^{-1}})_{ga}\circ\mathrm{d}(R_a)_g
		=\operatorname{Ad}_{a^{-1}}\circ\mathrm{d}(L_{g^{-1}})_g.
	\]
	因此右平移保持 \(h\) 恰好等价于所有 \(\operatorname{Ad}_{a^{-1}}\) 保持单位元内积；由于取逆遍历 \(G\)，这就是所述条件。紧群情形应用前一命题。

	若内积伴随不变，对 \(a=\exp_G(tZ)\) 在 \(t=0\) 求导，并用\cref{g1:prop:adjoint-differential-bracket}中已证的 \(\operatorname{ad}_ZX=[Z,X]\)，得到\eqref{g1:eq:ad-skew-inner-product}。反之，该式说明每个 \(\operatorname{ad}_Z\) 是斜伴随算子，于是
	\[
		\frac{\mathrm{d}}{\mathrm{d}t}
		\langle e^{t\operatorname{ad}_Z}X,e^{t\operatorname{ad}_Z}Y\rangle=0.
	\]
	\cref{g1:prop:exponential-naturality}给出 \(\operatorname{Ad}_{\exp_G(tZ)}=e^{t\operatorname{ad}_Z}\)，所以每个指数的伴随作用保持内积。\cref{g1:prop:exponentials-generate-identity-component}说明连通 \(G\) 的每个元素是有限个指数的乘积；正交变换的乘积仍正交，故得到整个 \(G\) 的伴随不变性。
\end{proof}

\begin{example}\label{g1:exm:unitary-trace-metric}
	在 \(\mathfrak u(n)=\{X:X^*=-X\}\) 上，取
	\[
		\langle X,Y\rangle=-\operatorname{Re}\operatorname{tr}(XY).
	\]
	它对称，且 \(\langle X,X\rangle=\operatorname{tr}(X^*X)>0\) 对 \(X\ne0\) 成立。由迹的循环性质，酉共轭保持它，故它给出 \(\operatorname{U}(n)\) 的双不变度量，也可限制到其 Lie 子群。并非每个左不变度量都双不变。例如在满足 \([H,E]=E\) 的二维仿射 Lie 代数上，\eqref{g1:eq:ad-skew-inner-product}若成立，就有 \(2\langle E,E\rangle=0\)，与正定性矛盾。
\end{example}

为了把一参数运动与“测地线”联系起来，还须指出这里用哪个联络。第15章已定义度量相容与挠率；双不变内积允许直接写出同时满足度量相容与无挠性的联络，再用它验证一参数子群的测地线方程。

\begin{proposition}[双不变度量的测地线]\label{g1:prop:bi-invariant-geodesics}
	设 \(G\) 为有限维实 Lie 群，配备双不变 Riemann 度量 \(h\)。则存在唯一保持 \(h\) 的无挠联络 \(\nabla\)，在左不变向量场上满足
	\begin{equation}\label{g1:eq:bi-invariant-connection}
		\nabla_{X^L}Y^L=\tfrac12[X,Y]^L.
	\end{equation}
	其测地线，即满足 \(\nabla_{\dot\gamma}\dot\gamma=0\) 的曲线，恰为
	\[
		\gamma(t)=g\exp_G(tX),\qquad g\in G,\ X\in\mathfrak g,
	\]
	或这些曲线的区间限制，参数起点可任意平移。
\end{proposition}
\begin{proof}
	取 \(\mathfrak g\) 的一组基，左平移给出整体标架 \(E_i\)。规定 \(\nabla_{E_i}E_j=\tfrac12[E_i,E_j]\)，再按第一输入的函数线性和第二输入的 Leibniz 法则延拓，便得到光滑联络；这是\cref{g1:prop:connection-local-data}中按标架指定联络系数的构造，此处标架覆盖全群，因而无需另行粘合。左不变场上的挠率为 \(\tfrac12[X,Y]-\tfrac12[Y,X]-[X,Y]=0\)，由\cref{g1:prop:torsion-tensorial}，挠率处处为零。左不变场的内积为常数，而\eqref{g1:eq:ad-skew-inner-product}给出
	\[
		h(\nabla_{Z^L}X^L,Y^L)+h(X^L,\nabla_{Z^L}Y^L)=0.
	\]
	用 Leibniz 法则展开一般场的分量，就得到度量相容条件。

	唯一性也可直接检查。两个这样的联络之差记为 \(B(U,V)\)；联络法则说明它是双线性张量。无挠性使 \(B(U,V)=B(V,U)\)，度量相容性使 \(b(U,V,W)=h(B(U,V),W)\) 关于后两个输入反对称。因此
	\[
		\begin{aligned}
			b(U,V,W)&=b(V,U,W)=-b(V,W,U)\\
			&=-b(W,V,U)=b(W,U,V)\\
			&=b(U,W,V)=-b(U,V,W).
		\end{aligned}
	\]
	故 \(b=0\)，正定性推出 \(B=0\)。这个唯一的无挠度量联络称为 Levi--Civita 联络。\footnote{Tullio Levi-Civita（1873--1941），意大利数学家，发展了绝对微分学与张量分析，对微分几何有重要贡献。}

	沿曲线写 \(\dot\gamma=\sum_i u_iE_i\)，并记左平移回单位元的速度为 \(u(t)=\sum_i u_i(t)E_i(e)\)。\cref{g1:def:along-curve-covariant-derivative}的沿曲线公式及\eqref{g1:eq:bi-invariant-connection}给出
	\[
		\nabla_{\dot\gamma}\dot\gamma
		=\sum_i u_i'E_i+\tfrac12\sum_{i,j}u_iu_j[E_i,E_j]
		=\sum_i u_i'E_i.
	\]
	最后一项因系数对称、括号反对称而抵消。所以测地线条件恰为 \(u'=0\)。设常值为 \(X\)，则曲线满足左不变场 \(X^L\) 的方程，\cref{g1:thm:lie-exponential}的存在唯一性给出 \(g\exp_G(tX)\)，并保证它对所有实参数有定义。
\end{proof}

\subsection{紧连通 Lie 群的最短测地线与指数满射}
\label{g1:subsec:1901:04}
测地线方程给出了每个初速度对应的曲线，却还没有说明它能否到达指定终点，也没有说明它是否最短。我们先证明短指数弧最短，再用紧性选择一条没有角点的最短折线。以下 \(G\) 配备双不变度量，\(|X|=\sqrt{\langle X,X\rangle}\) 表示单位元处的范数，\(|\dot c(t)|=\sqrt{h_{c(t)}(\dot c(t),\dot c(t))}\) 表示曲线的速度。分段 \(C^1\) 曲线 \(c:[a,b]\to G\) 的长度定义为
\[
	L(c)=\int_a^b|\dot c(t)|\,\mathrm{d}t,
\]
积分在各光滑段计算后相加。

\begin{lemma}[指数坐标中的径向估计]\label{g1:lem:compact-group-radial-estimate}
	设 \(G\) 为有限维实 Lie 群，具有双不变度量。在 \(\mathfrak g\) 中可取 \(r>0\)，使 \(\exp_G\) 在包含闭球 \(\overline B(0,2r)\) 的邻域上为到其像的微分同胚。记其局部逆为 \(\log\)。若 \(|X|<r\)，则从 \(e\) 到 \(\exp_GX\) 的任意分段 \(C^1\) 曲线长度至少为 \(|X|\)，而 \(t\mapsto\exp_G(tX)\)、\(0\le t\le1\)，的长度等于 \(|X|\)。从 \(e\) 到 \(G\setminus\exp_GB(0,r)\) 的任意此类曲线长度至少为 \(r\)。
\end{lemma}
\begin{proof}
	\cref{g1:prop:exponential-local-diffeomorphism}保证指数在零点附近可逆；取足够小的 \(r\) 即满足球的要求。\cref{g1:prop:lie-exponential-differential}说明，\((\mathrm{d}\exp_G)_ZY\) 左平移回单位元后等于
	\[
		J_ZY=\int_0^1e^{-s\operatorname{ad}_Z}Y\,\mathrm{d}s.
	\]
	伴随不变性使被积算子正交，而 \(e^{s\operatorname{ad}_Z}Z=Z\)，所以
	\begin{equation}\label{g1:eq:compact-group-radial-pairing}
		J_ZZ=Z,\qquad \langle J_ZY,Z\rangle=\langle Y,Z\rangle.
	\end{equation}
	这就是说，指数的微分保持径向速度，并把垂直于径向的向量仍送到垂直方向；它不声称所有方向的长度都保持。

	若曲线留在该指数坐标域中，写成 \(c(t)=\exp_GZ(t)\)，则速度范数为 \(|J_{Z(t)}Z'(t)|\)。为了允许曲线经过 \(e\)，对 \(\varepsilon>0\) 考虑光滑径向函数 \(r_\varepsilon(t)=\sqrt{|Z(t)|^2+\varepsilon^2}\)。由\eqref{g1:eq:compact-group-radial-pairing}与 Cauchy--Schwarz 不等式，
	\[
		|r_\varepsilon'(t)|
		=\frac{|\langle Z(t),J_{Z(t)}Z'(t)\rangle|}
		{\sqrt{|Z(t)|^2+\varepsilon^2}}
		\le |\dot c(t)|.
	\]
	在各光滑段积分、相加后令 \(\varepsilon\to0\)，得到
	\[
		\bigl||Z(b)|-|Z(a)|\bigr|\le L(c|_{[a,b]}).
	\]
	因此，从 \(e\) 出发而留在指数球内的曲线，其长度至少是终点的径向距离。

	还要排除先离开坐标域再绕回来的竞争曲线。若曲线从 \(e\) 首次离开 \(\exp_GB(0,r)\)，其首次到达的边界点形如 \(\exp_GZ\)、\(|Z|=r\)：这是因为闭球在较大的指数坐标域内，指数把闭球的边界准确映到该邻域的边界。前段曲线留在闭指数球内，上述估计给出其长度至少为 \(r\)。故当终点为 \(\exp_GX\)、\(|X|<r\) 时，任何离开指数球的竞争曲线都比 \(|X|\) 长；未离开的曲线则受径向估计控制。径向曲线的左平凡化速度恒为 \(X\)，长度正好 \(|X|\)，于是达到下界。
\end{proof}

这个局部估计也让我们能够严格使用全群的距离。设 \(G\) 连通，定义
\[
	d(p,q)=\inf_c L(c),
\]
其中 \(c:[a,b]\to G\) 遍历从 \(p\) 到 \(q\) 的分段 \(C^1\) 曲线。任意两点可由此类曲线连接：从一点可达的点集是开集，其补集也开，因为坐标小球内的直线像给出光滑连接；连通性排除了两个非空开集的分离。因此下确界有限。反向参数化与拼接分别给出对称性及三角不等式，左右平移保持长度，因而也保持 \(d\)。径向估计进一步给出
\begin{equation}\label{g1:eq:compact-group-local-distance}
	d(e,\exp_GX)=|X|\quad(|X|<r),\qquad
	B_d(e,r)=\exp_GB(0,r).
\end{equation}
其中第二式的反向包含用到了离开指数球至少花费长度 \(r\)。在更小半径重复同一估计，说明 \(d\) 正定且其开球给出原流形拓扑；左平移将结论搬到每一点。所以 \(d\) 确实是与原拓扑相容的距离。三角不等式还给出
\[
	|d(p,q)-d(p',q')|\le d(p,p')+d(q,q'),
\]
故 \(d\) 在 \(G\times G\) 上连续。

\begin{theorem}[紧群的最短测地线与指数满射]\label{g1:thm:compact-exp-surjective}
	设 \(G\) 为有限维实紧连通 Lie 群，并配备任一双不变 Riemann 度量。对任意 \(p,q\in G\)，存在 \(X\in\mathfrak g\)，使
	\[
		\gamma(t)=p\exp_G(tX),\qquad 0\le t\le1,
	\]
	连接 \(p,q\)，并且 \(L(\gamma)=|X|=d(p,q)\)。特别地，\(\exp_G:\mathfrak g\to G\) 满射。
\end{theorem}
\begin{proof}
	若 \(G\) 为单点，结论取 \(X=0\) 即得。以下设维数正。证明的难点是：一般的极小化曲线列可能没有足够好的导数收敛。我们只选择有限个中间点，把极小化化成紧空间上的连续函数问题，再利用局部指数弧消去所有角点。

	因为 \(d\) 连续而 \(G\times G\) 紧，其最大值 \(D\) 有限。取上一个引理中的 \(r>0\)，再取整数 \(N\ge2\) 满足 \(D/\sqrt N<r\)。固定 \(x_0=p\)、\(x_N=q\)，在紧空间 \(G^{N-1}\) 上定义
	\[
		E_N(x_1,\ldots,x_{N-1})
		=N\sum_{i=0}^{N-1}d(x_i,x_{i+1})^2.
	\]
	它连续，故有最小值 \(m\)。系数 \(N\) 对应把每一段安排在长度为 \(1/N\) 的参数区间上；最小化平方距离和，会同时惩罚过长的单段。

	我们先证 \(m\le d(p,q)^2\)。任取长度 \(L>0\) 的连接曲线，沿其累计弧长选取 \(N+1\) 个点，使相邻两点之间的弧长均为 \(L/N\)。累计弧长是从 \(0\) 到 \(L\) 的连续非减函数，因此即使原曲线有静止段，这些点也能按参数顺序选出。相邻两点的距离至多为 \(L/N\)，于是这个点组满足 \(E_N\le L^2\)。取长度趋于 \(d(p,q)\) 的曲线列，便有 \(m\le d(p,q)^2\)；若 \(p=q\)，也可直接取所有点等于 \(p\)。

	选定一个极小点组 \((x_i)\)。每段距离都满足
	\[
		d(x_i,x_{i+1})\le\sqrt{m/N}\le D/\sqrt N<r.
	\]
	所以\eqref{g1:eq:compact-group-local-distance}允许定义唯一的短向量
	\[
		X_i=\log(x_i^{-1}x_{i+1}),\qquad
		|X_i|=d(x_i,x_{i+1}).
	\]
	把短指数弧 \(x_i\exp_G(sX_i)\)、\(0\le s\le1\)，依次拼接，其长度为 \(\sum_i|X_i|\)。Cauchy--Schwarz 不等式及距离的定义给出
	\[
		d(p,q)\le\sum_i|X_i|
		\le\sqrt{N\sum_i|X_i|^2}=\sqrt m\le d(p,q).
	\]
	因此所有不等式都取等，折线已达到全局最短长度。但若只停在这一步，仍未说明折线在接点处光滑。

	为消去角点，先计算局部平方距离的一阶变化。在 \(\exp_GB(0,r)\) 上令 \(F(y)=\tfrac12|\log y|^2\)。若 \(y(t)=\exp_GZ(t)\)，其左平凡化速度为 \(v=\mathrm{d}(L_{y^{-1}})_y\dot y=J_ZZ'\)，由\eqref{g1:eq:compact-group-radial-pairing}有
	\begin{equation}\label{g1:eq:compact-group-distance-variation}
		\frac{\mathrm{d}}{\mathrm{d}t}F(y(t))
		=\langle Z,Z'\rangle=\langle Z,v\rangle.
	\end{equation}
	固定 \(1\le i\le N-1\)，把中间点改成 \(x_i(t)=x_i\exp_G(tV)\)，其中 \(V\in\mathfrak g\) 任意。其余点不动。由于相邻距离严格小于 \(r\)，小 \(t\) 时相关平方距离仍有光滑的上述表达。对前一段，左平移后为
	\[
		x_{i-1}^{-1}x_i(t)=\exp_G(X_{i-1})\exp_G(tV),
	\]
	在 \(t=0\) 的左平凡化速度为 \(V\)，所以半平方距离的导数为 \(\langle X_{i-1},V\rangle\)。对后一段先用距离的对称性，再左平移成
	\[
		x_{i+1}^{-1}x_i(t)=\exp_G(-X_i)\exp_G(tV).
	\]
	其半平方距离导数为 \(-\langle X_i,V\rangle\)。极小点的一阶导数为零，因此
	\[
		0=\left.\frac{\mathrm{d}}{\mathrm{d}t}\right|_{t=0}E_N
		=2N\langle X_{i-1}-X_i,V\rangle
		\qquad(V\in\mathfrak g).
	\]
	内积非退化，故 \(X_{i-1}=X_i\)。所有短向量于是等于同一个 \(Y\)，从而 \(x_i=p\exp_G(iY)\)。此前的折线以每段相同的时间 \(1/N\) 参数化，恰好成为一条光滑曲线 \(p\exp_G(tNY)\)。令 \(X=NY\)，由\cref{g1:prop:bi-invariant-geodesics}它是测地线，而已经证明的长度等式给出 \(|X|=d(p,q)\)。最后取 \(p=e\)、\(q\) 任意，便有 \(q=\exp_GX\)。
\end{proof}

\begin{remark}\label{g1:rem:compact-exponential-boundaries}
	紧性在证明中提供有限直径和有限点组极小值，双不变性则提供径向配对及相接速度的统一表达。结论不包括指数单射，也不说每条测地线段都最短：在标准圆周群上，\(t\mapsto e^{\mathrm{i}t\theta}\)、\(0\le t\le1\)，若 \(|\theta|>\pi\)，便可换一条更短的圆弧。若紧群不连通，指数的像只能落在 \(G^0\)，因为 \(t\mapsto\exp_G(tX)\) 把每个指数连到 \(e\)；上述定理应用于紧连通的 \(G^0\)，说明指数的像恰为 \(G^0\)。
\end{remark}
% 纲要条目排为编号小节标题，并进入目录与书签；说明段仍保留于纲要注释中。

% !TeX root = ../../Book1.tex
% 纲要标题：### 19.2 极大环面的存在
% 纲要位置：几何学计划纲要.md，第 853 行。
% 正文状态：初稿；尚未正式审核。

% 纲要内容开始
% 
% * 一参数子群的闭包
% * 紧连通交换 Lie 群的结构
% * 极大环面的构造
% * 每个元素包含在极大环面中的证明
% 
% 纲要内容结束

\section{极大环面的存在}
\label{g1:sec:1902}

圆周群中的运动可以反复绕行，两个圆周的乘积还允许以互不成比例的速度同时转动。后一种运动虽然只有一个实参数，其像的闭包却可能是整个二维环面。因此，在紧群中寻找交换方向时，合适的对象不是单参数子群的像本身，而是它的闭包。闭子群定理赋予这个闭包 Lie 群结构；紧连通交换 Lie 群的结构定理则将说明，它正是若干圆周的乘积。

\subsection{一参数子群的闭包}
\label{g1:subsec:1902:01}
设 \(G\) 为紧 Lie 群，\(X\in\mathfrak g\)。记
\[
	A_X\coloneqq\overline{\{\exp_G(tX):t\in\mathbb R\}},
\]
其中闭包取在 \(G\) 中。参数化的像可能不是嵌入子流形；其闭包却是闭子群，可以使用闭子群定理。连通性与交换性也传到闭包，因而将得到一个环面。

\begin{lemma}[交换子群的闭包]\label{g1:lem:abelian-subgroup-closure}
	设 \(G\) 为有限维实 Lie 群，\(A\subseteq G\) 为交换子群。则 \(\overline A\) 仍是交换子群。若 \(A\) 连通，则 \(\overline A\) 连通；若 \(G\) 紧，则 \(\overline A\) 是紧 Lie 子群。特别地，对 \(X\in\operatorname{Lie}(G)\)，令 \(A_X=\overline{\{\exp_G(tX):t\in\mathbb R\}}\)；若 \(G\) 紧，则 \(A_X\) 是紧连通交换 Lie 群。
\end{lemma}
\begin{proof}
	Lie 群为局部 Euclid 且第二可数的空间，点的闭包关系可以用序列检测。取 \(a,b\in\overline A\)，选 \(a_j,b_j\in A\) 分别趋于它们。连续性给出 \(a_jb_j^{-1}\to ab^{-1}\)，所以 \(ab^{-1}\in\overline A\)，这证明闭包为子群。又因 \(a_jb_j=b_ja_j\)，取极限得 \(ab=ba\)。

	连通集的闭包连通：若闭包被两个非空相对开闭集分开，稠密子集 \(A\) 必须与两者都相交，于是也被分开，矛盾。紧群的闭子集紧，而\cref{g1:thm:closed-subgroup}使闭子群成为采用子空间拓扑的嵌入 Lie 子群。最后，\(t\mapsto\exp_G(tX)\) 是连通空间 \(\mathbb R\) 的连续像，且单参数乘法规律使像为交换子群，因此适用于上述结论。
\end{proof}

\cref{g1:exm:irrational-torus-subgroup}中的 \(t\mapsto(e^{\mathrm{i}t},e^{\mathrm{i}\alpha t})\)、\(\alpha\notin\mathbb Q\)，在 \(S^1\times S^1\) 中稠密。这说明 \(\operatorname{Lie}(A_X)\) 可能严格包含 \(\mathbb RX\)：闭包加入了极限方向，不能按最初参数的个数计算闭包的维数。相反，若这条单参数子群有非零周期，它的像已是一个圆周，因而无需再添极限点。\(X=0\) 时闭包只是单位元。

\subsection{紧连通交换 Lie 群的结构}
\label{g1:subsec:1902:02}
现在从一个抽象的紧连通交换 Lie 群 \(A\) 出发。其 Lie 代数 \(\mathfrak a\) 的括号为零，故\cref{g1:cor:commuting-exponentials}表明指数映射保持加法：
\[
	\exp_A(X+Y)=\exp_A X\,\exp_A Y.
\]
指数在零点附近是微分同胚，所以像是包含单位元邻域的子群，因而是开子群。\(A\) 连通，开子群的各陪集又都是开集，故像只能为整个 \(A\)。因此分类问题变成：指数映射的核在向量空间里是什么样的子群？

\begin{lemma}[离散子群的格基]\label{g1:lem:discrete-subgroup-lattice-basis}
	设 \(V\) 为 \(r\) 维实向量空间，\(\Lambda\subset V\) 为离散加法子群。则存在实线性无关的向量 \(v_1,\ldots,v_k\)，其中 \(0\le k\le r\)，使
	\[
		\Lambda=\mathbb Zv_1\oplus\cdots\oplus\mathbb Zv_k.
	\]
	商空间 \(V/\Lambda\) 紧当且仅当 \(k=r\)。当 \(k=r\) 时，称 \(\Lambda\) 为 \(V\) 中的\term[manzhige]{满秩格}[full lattice]。
\end{lemma}
\begin{proof}
	在 \(V\) 上选 Euclid 内积。离散性给出 \(\varepsilon>0\)，使 \(\Lambda\cap B_\varepsilon(0)=\{0\}\)。不同格点之差仍在 \(\Lambda\) 中，所以不同格点的距离至少为 \(\varepsilon\)。这也说明每个有界集只包含有限多个格点，并且 \(\Lambda\) 闭：收敛的格点序列为 Cauchy 序列，最终必须恒定。

	对 \(r\) 归纳。\(\Lambda=\{0\}\) 时取空基。否则选长度最小的非零 \(v_1\in\Lambda\)；最小值存在，因为先取一个非零格点，其长度以内的格点只有有限个。若 \(sv_1\in\Lambda\)，减去 \(\lfloor s\rfloor v_1\) 后得到长度小于 \(\lVert v_1\rVert\) 的格点，故余项为零。这证明
	\[
		\Lambda\cap\mathbb Rv_1=\mathbb Zv_1.
	\]
	令 \(p:V\to v_1^\perp\) 为正交投影。我们需要先证明 \(p(\Lambda)\) 仍离散；任意离散集合的投影并没有这一性质。若存在非零 \(p(\lambda_j)\to0\)，将 \(\lambda_j\) 减去适当整数倍的 \(v_1\)，可使其沿 \(v_1\) 的分量都位于线段 \([0,v_1]\) 内。这些调整后的格点有界，因而只有有限种；其非零投影便不可能趋于零，矛盾。

	由归纳假设，\(p(\Lambda)\) 有格基 \(w_2,\ldots,w_k\)。各取一个提升 \(v_j\in\Lambda\)，满足 \(p(v_j)=w_j\)。任意 \(\lambda\in\Lambda\) 的投影都是这些 \(w_j\) 的整数线性组合，减去相应 \(v_j\) 后，余项位于 \(\Lambda\cap\mathbb Rv_1=\mathbb Zv_1\)。因此 \(v_1,\ldots,v_k\) 生成 \(\Lambda\)。若它们有实线性关系，先投影便得到 \(w_2,\ldots,w_k\) 的实线性关系，迫使相应系数为零，再由 \(v_1\ne0\) 得剩余系数为零。

	记 \(W=\operatorname{span}_{\mathbb R}\Lambda\)。映射 \(V/\Lambda\to V/W\)、\([v]\mapsto v+W\) 连续且满射。若 \(V/\Lambda\) 紧，则向量空间 \(V/W\) 紧，只能有 \(W=V\)，即 \(k=r\)。反之，\(k=r\) 时每个陪集都有代表元位于紧平行多面体
	\[
		\left\{\sum_{j=1}^r t_jv_j:0\le t_j\le1\right\},
	\]
	因为可将各坐标减去整数部分。商空间是这个紧集的连续像，所以紧。
\end{proof}

\begin{theorem}[紧连通交换 Lie 群的结构]\label{g1:thm:compact-connected-abelian-group}
	设 \(A\) 为 \(r\) 维紧连通交换实 Lie 群，\(\mathfrak a=\operatorname{Lie}(A)\)。则
	\[
		\Lambda_A\coloneqq\ker(\exp_A:\mathfrak a\to A)
	\]
	是 \(\mathfrak a\) 中的满秩格，指数诱导 Lie 群同构
	\[
		\mathfrak a/\Lambda_A\ \cong\ A\ \cong\ (S^1)^r.
	\]
	\(r=0\) 时右侧解释为平凡群。
\end{theorem}
\begin{proof}
	前面的论证已给出指数为满的光滑同态。由\cref{g1:prop:exponential-local-diffeomorphism}，零点有一个邻域只包含核的零元，故 \(\Lambda_A\) 离散；它也是闭集 \(\{e\}\) 的原像。\cref{g1:thm:homogeneous-space}应用于加法 Lie 群 \(\mathfrak a\) 及其闭子群 \(\Lambda_A\)，赋予商空间光滑结构。指数在陪集上常值，因此下降成光滑双射 \(\overline{\exp}_A:\mathfrak a/\Lambda_A\to A\)。

	这不仅是一个集合双射。商映射与指数映射在零点附近都是局部微分同胚，故 \(\overline{\exp}_A\) 在单位元附近具有光滑局部逆。它又是同态，用群平移即可将局部逆移到每一点。因此它是 Lie 群同构，商空间也由 \(A\) 的紧性而紧。\cref{g1:lem:discrete-subgroup-lattice-basis}于是给出满秩格基 \(v_1,\ldots,v_r\)。线性坐标 \(\sum_jt_jv_j\mapsto(t_1,\ldots,t_r)\) 将格化为 \(\mathbb Z^r\)，再用
	\[
		(t_1,\ldots,t_r)+\mathbb Z^r\longmapsto
		(e^{2\pi\mathrm{i}t_1},\ldots,e^{2\pi\mathrm{i}t_r})
	\]
	得到到 \((S^1)^r\) 的 Lie 群同构。
\end{proof}

选择格基才得到与标准圆周积的具体同构；指数核本身不依赖这项选择。以后比较具有相同 Lie 代数的紧群时，核格的差别将记录部分整体信息。这里的 \(2\pi\) 来自圆周的角度坐标：若用 \(\mathrm{i}\mathbb R\) 表示 \(\operatorname{Lie}(S^1)\)，指数核为 \(2\pi\mathrm{i}\mathbb Z\)，而不是 \(\mathrm{i}\mathbb Z\)。

\subsection{极大环面的构造}
\label{g1:subsec:1902:03}
\begin{definition}\label{g1:def:maximal-torus}
	设 \(G\) 为有限维实 Lie 群。\(G\) 中的紧连通交换 Lie 子群称为\term[huanmian]{环面}[torus]。若一个环面 \(T\subseteq G\) 不真包含于任何更大的环面，则称 \(T\) 为\term[jidahuanmian]{极大环面}[maximal torus]。平凡群也允许作为零维环面。\symidx{\(T,\mathfrak t\)：极大环面及其 Lie 代数}
\end{definition}

这个定义中的“极大”按包含关系判断。它并不表示存在一个包含全部环面的“最大环面”。例如绕不同轴旋转得到 \(\operatorname{SO}(3)\) 中不同的圆周子群，它们彼此不包含。随后证明的共轭定理将说明，这些极大环面仍具有共同的维数。

\begin{proposition}[环面的极大扩张]\label{g1:prop:maximal-torus-existence}
	设 \(G\) 为有限维紧实 Lie 群。任意环面 \(T_0\subseteq G\) 都包含于一个极大环面。特别地，\(G\) 存在极大环面。
\end{proposition}
\begin{proof}
	考虑全部包含 \(T_0\) 的环面，其维数是 \(0,1,\ldots,\dim G\) 中的整数，且这个集合非空。取其中维数最大的一个 \(T\)。若 \(T\subseteq S\) 而 \(S\) 也是环面，维数最大性给出 \(\dim T=\dim S\)。包含映射在单位元处的微分因维数相同而可逆，故 \(T\) 包含 \(S\) 中单位元的一个开邻域。于是 \(T\) 是连通群 \(S\) 的开子群，只能有 \(T=S\)。这证明 \(T\) 极大。取 \(T_0=\{e\}\) 即得存在性。
\end{proof}

相同论证说明，环面之间的真包含必使维数严格增加；因此此处的极大扩张只需有限维数论证。\cref{g1:lem:abelian-subgroup-closure,g1:thm:compact-connected-abelian-group}又表明，每条单参数子群的闭包 \(A_X\) 是一个环面，所以任意 \(X\in\mathfrak g\) 都属于某个极大环面的 Lie 代数：包含 \(A_X\) 的环面包含整条 \(\exp_G(tX)\)，在 \(t=0\) 求导即可。

\subsection{每个元素包含在极大环面中的证明}
\label{g1:subsec:1902:04}
上一步已经处理所有无穷小方向。要把任意群元素也放入环面，关键是它能否写成一个指数。\cref{g1:thm:compact-exp-surjective}对紧连通群给出了肯定答案；连通性在这里真正进入。

\begin{theorem}[元素属于极大环面]\label{g1:thm:element-in-maximal-torus}
	设 \(G\) 为紧连通实 Lie 群。对每个 \(g\in G\)，存在极大环面 \(T\subseteq G\)，使 \(g\in T\)。
\end{theorem}
\begin{proof}
	由紧连通群的指数满射性，取 \(X\in\mathfrak g\) 满足 \(g=\exp_GX\)。单参数子群的闭包 \(A_X\) 是一个环面，并且包含 \(g\)。由\cref{g1:prop:maximal-torus-existence}，它可扩张到极大环面 \(T\)，于是 \(g\in A_X\subseteq T\)。
\end{proof}

若 \(G\) 不连通，任何环面由于连通且包含单位元，都位于 \(G^0\)，所以别的连通分支中的元素不可能包含在环面中。例如 \(\operatorname{O}(2)\) 中的反射不在任何环面中。对紧连通群，定理则给出
\[
	G=\bigcup_{T\text{ 为极大环面}}T.
\]
这个等式暂时允许不同环面具有不同位置。下一节将证明，它们全部通过群内共轭相互得到，于是可以用一个固定环面研究所有群元素的共轭类。

% !TeX root = ../../Book1.tex
% 纲要标题：### 19.3 共轭性与中心化子
% 纲要位置：几何学计划纲要.md，第 860 行。
% 正文状态：初稿；尚未正式审核。

% 纲要内容开始
% 
% * 环面作用及无穷小分解
% * 正则元素和中心化子
% * 极大环面的共轭定理
% * 环面内共轭与正规化子
% 
% 纲要内容结束

\section{共轭性与中心化子}
\label{g1:sec:1903}

极大环面已经存在，但目前它可能随群元素改变。为了比较两个环面，先研究一个固定环面在 Lie 代数上的伴随作用。交换性使这些线性变换可以同时分解，紧性使它们成为旋转；由此选出一个向量，使与它交换的方向尽量少。将任意伴随轨道与这个向量作内积，再利用紧性取极值，就能把整条轨道带到固定环面的 Lie 代数中。

\ProseParagraphBreak

以下 \(G\) 为紧连通实 Lie 群，\(T\subseteq G\) 为一个极大环面，\(\mathfrak t=\operatorname{Lie}(T)\)。由\cref{g1:prop:bi-invariant-metric}，\(\mathfrak g\) 上有正定的 \(\operatorname{Ad}(G)\) 不变内积 \(\langle\cdot,\cdot\rangle\)。本节每次取正交补或内积极值，都使用这同一个内积。

\subsection{环面作用及无穷小分解}
\label{g1:subsec:1903:01}
对 \(H\in\mathfrak t\)，算子 \(\operatorname{ad}_H:X\mapsto[H,X]\) 是斜伴随的。事实上，\(\operatorname{Ad}_{\exp(sH)}\) 保持内积，对
\[
	\langle\operatorname{Ad}_{\exp(sH)}X,
	\operatorname{Ad}_{\exp(sH)}Y\rangle=\langle X,Y\rangle
\]
在 \(s=0\) 求导，得到
\begin{equation}\label{g1:eq:invariant-inner-product-bracket}
	\langle[H,X],Y\rangle+\langle X,[H,Y]\rangle=0.
\end{equation}
由于 \(\mathfrak t\) 交换，Jacobi 恒等式还给出 \([\operatorname{ad}_H,\operatorname{ad}_K]=\operatorname{ad}_{[H,K]}=0\)。因此我们面对的是一族彼此交换的斜伴随算子。

\begin{lemma}[交换斜 Hermitian 算子的同时分解]\label{g1:lem:commuting-skew-hermitian}
	设 \(V\) 为有限维复 Hermitian 空间，\(A_1,\ldots,A_r\) 为彼此交换的斜 Hermitian 线性算子，即 \(A_j^*=-A_j\)。则 \(V\) 有一组共同特征向量组成的标准正交基；各 \(A_j\) 的特征值均为纯虚数。
\end{lemma}
\begin{proof}
	先看一个算子 \(A\)。复数域上的特征多项式有根，故有单位特征向量 \(v\)。若 \(Av=\lambda v\)，斜 Hermitian 性给出 \(\lambda=-\overline\lambda\)，所以 \(\lambda\) 纯虚。\(v^\perp\) 在 \(A\) 下不变，因为对 \(w\perp v\)，有 \(\langle Aw,v\rangle=\langle w,-Av\rangle=0\)。在 \(v^\perp\) 上归纳，得到标准正交特征基。

	现在先分解 \(A_1\) 的各特征空间。由 \(A_jA_1=A_1A_j\)，每个 \(A_j\) 保持这些空间，其限制仍斜 Hermitian。在每个空间内再按 \(A_2\) 分解，继续至 \(A_r\)。每一步只细分已有的正交分解，所以最后得到共同特征向量的标准正交基。
\end{proof}

为使用此引理，把 \(\mathfrak g\) 复化为 \(\mathfrak g_{\mathbb C}=\mathfrak g\otimes_{\mathbb R}\mathbb C\)。原实内积复双线性延拓为 \(B_{\mathbb C}\)，再定义 Hermitian 内积 \((Z,W)=B_{\mathbb C}(Z,\overline W)\)，其中共轭固定实空间 \(\mathfrak g\)，本式对第一变量线性。\(\operatorname{ad}_H\) 复线性延拓后是斜 Hermitian 算子。取 \(\mathfrak t\) 的一组基同时对角化，每个共同特征空间上的特征值随 \(H\) 实线性变化，可以唯一写成 \(\mathrm{i}\alpha(H)\)，其中 \(\alpha\in\mathfrak t^*\) 为实线性泛函。

\begin{proposition}[环面伴随作用的权分解]\label{g1:prop:torus-weight-decomposition}
	设 \(G\) 为紧连通实 Lie 群，\(T\subseteq G\) 为环面，\(\mathfrak g=\operatorname{Lie}(G)\)、\(\mathfrak t=\operatorname{Lie}(T)\)。固定 \(\mathfrak g\) 上的 \(\operatorname{Ad}(G)\) 不变正定实内积 \(B\)，在 \(\mathfrak g_{\mathbb C}\) 上取其 Hermitian 延拓，约定第一变量线性、第二变量共轭线性。存在有限集 \(R\subseteq\mathfrak t^*\setminus\{0\}\)，使
	\begin{equation}\label{g1:eq:torus-adjoint-weight-decomposition}
		\mathfrak g_{\mathbb C}=\mathfrak g_0\oplus
		\bigoplus_{\alpha\in R}\mathfrak g_\alpha,
		\qquad
		\mathfrak g_\alpha=\{Z:[H,Z]=\mathrm{i}\alpha(H)Z
		\;\text{对所有 }H\in\mathfrak t\}.
	\end{equation}
	这些空间相对于复化后所得的 Hermitian 内积彼此正交，且 \(\overline{\mathfrak g_\alpha}=\mathfrak g_{-\alpha}\)。零权空间为
	\[
		\mathfrak g_0=\{X\in\mathfrak g:[X,H]=0
		\;\text{对所有 }H\in\mathfrak t\}_{\mathbb C}.
	\]
	若 \(\Lambda_T=\ker(\exp_G|_{\mathfrak t})\)，则对 \(\alpha\in R\) 有 \(\alpha(\Lambda_T)\subseteq2\pi\mathbb Z\)，并且
	\begin{equation}\label{g1:eq:torus-weight-character}
		\operatorname{Ad}_{\exp_GH}Z=e^{\mathrm{i}\alpha(H)}Z
		\qquad(Z\in\mathfrak g_\alpha).
	\end{equation}
\end{proposition}
\begin{proof}
	前面的同时分解给出有限个实线性泛函。合并相同泛函的共同特征空间，并将非零者组成 \(R\)，就得到直和分解。若两个泛函不同，可取某个 \(H\) 使其取值不同；斜 Hermitian 算子 \(\operatorname{ad}_H\) 的不同特征空间正交，所以整个分解正交。零权条件就是与 \(\mathfrak t\) 的所有元素交换；将 \(Z=X+\mathrm{i}Y\) 分成实虚部可知它是所列实空间的复化。

	由于 \(H\) 为实向量，对 \([H,Z]=\mathrm{i}\alpha(H)Z\) 取共轭得到 \([H,\overline Z]=-\mathrm{i}\alpha(H)\overline Z\)，所以共轭交换 \(\mathfrak g_\alpha\) 与 \(\mathfrak g_{-\alpha}\)。由\cref{g1:prop:exponential-naturality}，\(\operatorname{Ad}_{\exp_GH}=e^{\operatorname{ad}_H}\)。将算子指数作用在特征空间上，便得\eqref{g1:eq:torus-weight-character}。若 \(H\in\Lambda_T\)，左侧为恒等作用；\(\mathfrak g_\alpha\ne0\) 迫使 \(e^{\mathrm{i}\alpha(H)}=1\)，即 \(\alpha(H)\in2\pi\mathbb Z\)。
\end{proof}

我们把 \(\alpha\) 称为这一作用的\term[wuqiongxiaoquan]{无穷小权}[infinitesimal weight]。其群上版本是\term[tezhengbiao]{特征标}[character]，即光滑群同态 \(\chi_\alpha:T\to S^1\)，这里由 \(\chi_\alpha(\exp H)=e^{\mathrm{i}\alpha(H)}\) 给出。\cref{g1:thm:compact-connected-abelian-group}保证每个 \(t\in T\) 都有这样的指数表示，而刚证明的整性保证换一个相差格向量的 \(H\) 不改变取值。复化将实旋转写成标量乘法；共轭配对说明这些复特征空间仍来自实几何。

\subsection{正则元素和中心化子}
\label{g1:subsec:1903:02}
对一个群子集 \(A\subseteq G\)，与它的每个元素都交换的群元素组成\term[zhongxinhua]{中心化子}[centralizer]
\[
	Z_G(A)\coloneqq\{g\in G:ga=ag\;\text{对所有 }a\in A\}.
\]
它是闭子群，因为每个交换条件都是连续映射的闭水平集。由\cref{g1:thm:closed-subgroup}，其 Lie 代数可以用全部留在中心化子中的单参数子群检测；指数自然性给出
\[
	\operatorname{Lie}(Z_G(A))
	=\{X\in\mathfrak g:\operatorname{Ad}_aX=X\;\text{对所有 }a\in A\}.
\]
这里两个方向都可直接核对：若 \(\exp(sX)\) 与 \(a\) 交换，微分得 \(\operatorname{Ad}_aX=X\)；反之此等式通过指数自然性恢复全部 \(s\) 的交换关系。对于 Lie 代数子集 \(E\subseteq\mathfrak g\)，相应记号为
\[
	\mathfrak z_{\mathfrak g}(E)\coloneqq\{X:[X,Y]=0\;\text{对所有 }Y\in E\}.
\]

\begin{proposition}[极大环面的无穷小中心化子]\label{g1:prop:maximal-torus-lie-centralizer}
	设 \(G\) 为紧连通实 Lie 群，\(T\subseteq G\) 为极大环面，\(\mathfrak t=\operatorname{Lie}(T)\)。则
	\[
		\operatorname{Lie}(Z_G(T))=\mathfrak z_{\mathfrak g}(\mathfrak t)=\mathfrak t.
	\]
	特别地，\(\mathfrak t\) 是 \(\mathfrak g\) 的极大交换 Lie 子代数，且权分解的零权空间为 \(\mathfrak g_0=\mathfrak t_{\mathbb C}\)。
\end{proposition}
\begin{proof}
	因为 \(T=\exp_G(\mathfrak t)\)，与全部 \(\operatorname{Ad}_t\)、\(t\in T\)，相容的不动条件等价于 \([H,X]=0\) 对所有 \(H\in\mathfrak t\) 成立：一个方向对 \(t=\exp(sH)\) 求导，另一个方向用 \(\operatorname{Ad}_{\exp H}=e^{\operatorname{ad}_H}\)。这证明前两个空间相等。

	若 \(X\in\mathfrak z_{\mathfrak g}(\mathfrak t)\)，则 \(\mathfrak a=\mathfrak t+\mathbb RX\) 是交换 Lie 子代数。由\cref{g1:cor:commuting-exponentials}，\(\exp_G(\mathfrak a)\) 是连通交换子群。它的闭包由\cref{g1:lem:abelian-subgroup-closure}是紧连通交换 Lie 群，因而是包含 \(T\) 的环面。\(T\) 的极大性迫使这个闭包等于 \(T\)。于是 \(\exp_G(sX)\in T\) 对全部 \(s\) 成立，在零点求导得 \(X\in\mathfrak t\)。反向包含来自 \(\mathfrak t\) 本身交换。任何包含 \(\mathfrak t\) 的交换 Lie 子代数都在其中心化子中，所以只能等于 \(\mathfrak t\)。零权空间的结论随即由\eqref{g1:eq:torus-adjoint-weight-decomposition}得到。
\end{proof}

对 \(H\in\mathfrak t\)，权分解现在给出更精确的式子
\begin{equation}\label{g1:eq:regular-element-centralizer}
	\mathfrak z_{\mathfrak g}(H)_{\mathbb C}
	=\mathfrak t_{\mathbb C}\oplus
	\bigoplus_{\substack{\alpha\in R\\\alpha(H)=0}}\mathfrak g_\alpha.
\end{equation}
所以，想使与 \(H\) 交换的方向尽量少，只需让 \(H\) 避开有限多个超平面 \(\ker\alpha\)。

\begin{definition}\label{g1:def:regular-toral-element}
	设 \(T\) 为紧连通 Lie 群 \(G\) 的极大环面。满足 \(\mathfrak z_{\mathfrak g}(H)=\mathfrak t\) 的 \(H\in\mathfrak t\) 称为\term[zhengzeyuansu]{正则元素}[regular element]。它们组成的集合记为 \(\mathfrak t_{\mathrm{reg}}\)。
\end{definition}

\begin{lemma}[正则元素的存在]\label{g1:lem:regular-toral-elements}
	设 \(G\) 为紧连通实 Lie 群，\(T\) 为其极大环面，\(\mathfrak g=\operatorname{Lie}(G)\)、\(\mathfrak t=\operatorname{Lie}(T)\)。令 \(R\) 为所有满足下列条件的非零实泛函 \(\alpha\in\mathfrak t^*\) 的集合：存在非零 \(Z\in\mathfrak g_{\mathbb C}\)，使 \([H,Z]=\mathrm{i}\alpha(H)Z\) 对每个 \(H\in\mathfrak t\) 成立。以 \(\mathfrak t_{\mathrm{reg}}\) 记在 \(\mathfrak g\) 中中心化子恰为 \(\mathfrak t\) 的 \(H\in\mathfrak t\) 所成的集合。则
	\[
		\mathfrak t_{\mathrm{reg}}=
		\mathfrak t\setminus\bigcup_{\alpha\in R}\ker\alpha
	\]
	是 \(\mathfrak t\) 中非空的稠密开集。若 \(R=\varnothing\)，则 \(\mathfrak t_{\mathrm{reg}}=\mathfrak t\)。
\end{lemma}
\begin{proof}
	所列等式来自\eqref{g1:eq:regular-element-centralizer}。每个 \(\alpha\ne0\) 的核都是闭的真超平面，其补集开且稠密：在任意开球中沿一个 \(\alpha\) 不为零的方向稍作移动，就能离开超平面。有限多个稠密开集的交仍稠密，可依次在给定非空开集中选一个更小的非空开球避开下一个超平面。这个交也是开集，故有结论。空根集时没有需要避开的超平面。
\end{proof}

这里的正则性首先属于 Lie 代数。对群元素 \(\exp H\)，中心化子的 Lie 代数却由 \(\operatorname{Ad}_{\exp H}\) 的不动向量给出，条件是 \(\alpha(H)\in2\pi\mathbb Z\)，而非 \(\alpha(H)=0\)。例如在 \(\operatorname{SU}(2)\) 中，\(H=\operatorname{diag}(\mathrm{i}\pi,-\mathrm{i}\pi)\) 的 Lie 代数中心化子是对角子代数，但 \(\exp H=-I\) 与全群交换。这一差别将在处理群元素的共轭时再次出现。

\subsection{极大环面的共轭定理}
\label{g1:subsec:1903:03}
固定一个正则 \(H\in\mathfrak t\)。若某个向量 \(Z\) 与 \(H\) 交换，正则性就迫使 \(Z\in\mathfrak t\)。因此，把任意 \(X\in\mathfrak g\) 共轭入 \(\mathfrak t\)，可以转化为寻找一个共轭代表，使它与 \(H\) 交换。内积函数在极值点的一阶条件恰好给出这个括号方程。

\begin{proposition}[伴随轨道与固定环面相交]\label{g1:prop:adjoint-orbit-meets-torus}
	设 \(G\) 为紧连通实 Lie 群，\(T\subseteq G\) 为一个极大环面。对任意 \(X\in\mathfrak g\)，存在 \(g\in G\) 使 \(\operatorname{Ad}_gX\in\mathfrak t\)。
\end{proposition}
\begin{proof}
	选定 \(\operatorname{Ad}(G)\) 不变正定内积，以及\cref{g1:lem:regular-toral-elements}给出的 \(H\in\mathfrak t_{\mathrm{reg}}\)。连续函数
	\[
		f:G\to\mathbb R,\qquad f(g)=\langle\operatorname{Ad}_gX,H\rangle
	\]
	在紧空间 \(G\) 上达到最大值。设最大点为 \(g_0\)，记 \(Z=\operatorname{Ad}_{g_0}X\)。对任意 \(Y\in\mathfrak g\)，曲线 \(s\mapsto\exp(sY)g_0\) 经过最大点，故
	\[
		0=\left.\frac{\mathrm{d}}{\mathrm{d}s}\right|_{s=0}
		f(\exp(sY)g_0)
		=\langle[Y,Z],H\rangle
		=\langle Y,[Z,H]\rangle.
	\]
	最后一等式由\eqref{g1:eq:invariant-inner-product-bracket}及括号反对称性得到。它对所有 \(Y\) 成立，特别取 \(Y=[Z,H]\)，正定性给出 \([Z,H]=0\)。正则性说明 \(Z\in\mathfrak z_{\mathfrak g}(H)=\mathfrak t\)，所以 \(g=g_0\) 满足要求。
\end{proof}

\begin{theorem}[极大环面的共轭定理]\label{g1:thm:maximal-tori-conjugate}
	设 \(G\) 为紧连通实 Lie 群。任意两个极大环面 \(T,S\subseteq G\) 共轭，即存在 \(g\in G\)，使 \(gSg^{-1}=T\)。因此所有极大环面的维数相等。
\end{theorem}
\begin{proof}
	记 \(\mathfrak t=\operatorname{Lie}(T)\)、\(\mathfrak s=\operatorname{Lie}(S)\)。在 \(\mathfrak s\) 中选择正则元素 \(X\)，使 \(\mathfrak z_{\mathfrak g}(X)=\mathfrak s\)。由\cref{g1:prop:adjoint-orbit-meets-torus}，有 \(g\in G\) 使 \(\operatorname{Ad}_gX\in\mathfrak t\)。由于 \(\mathfrak t\) 交换，得到
	\[
		\mathfrak t\subseteq\mathfrak z_{\mathfrak g}(\operatorname{Ad}_gX)
		=\operatorname{Ad}_g\mathfrak z_{\mathfrak g}(X)
		=\operatorname{Ad}_g\mathfrak s.
	\]
	中间的等式来自共轭保持括号。右端是交换 Lie 子代数，而\cref{g1:prop:maximal-torus-lie-centralizer}说明 \(\mathfrak t\) 已经是极大交换 Lie 子代数，所以这个包含只能是等号。

	两个环面都是连通 Lie 子群，且 \(gSg^{-1}\) 与 \(T\) 的 Lie 代数相同；\cref{g1:thm:lie-subalgebra-integration}中的连通积分子群唯一性给出 \(gSg^{-1}=T\)。也可直接用环面指数的满射性以及指数与共轭相容，写成 \(gSg^{-1}=\exp_G(\operatorname{Ad}_g\mathfrak s)=\exp_G\mathfrak t=T\)。共轭是微分同胚，故维数相同。
\end{proof}

\begin{definition}\label{g1:def:compact-group-rank}
	紧连通 Lie 群 \(G\) 的任一极大环面的维数称为 \(G\) 的\term[jinqunzhizhishu]{秩}[rank]，记为 \(\operatorname{rank}G\)。
\end{definition}

秩的选择无关性现在来自共轭定理。由\cref{g1:prop:adjoint-orbit-meets-torus}及\eqref{g1:eq:regular-element-centralizer}，对任意 \(X\in\mathfrak g\) 都有 \(\dim\mathfrak z_{\mathfrak g}(X)\ge\operatorname{rank}G\)，而正则元素恰使这个维数取最小值。因此可以对一般 \(X\in\mathfrak g\) 用最小中心化子维数定义正则性，与刚才在 \(\mathfrak t\) 中的定义一致。

\begin{proposition}[极大环面的中心化子]\label{g1:prop:maximal-torus-centralizer}
	设 \(G\) 为紧连通实 Lie 群，\(T\subseteq G\) 为极大环面。则 \(Z_G(T)=T\)。
\end{proposition}
\begin{proof}
	\(T\subseteq Z_G(T)\) 来自交换性。反过来，取 \(a\in Z_G(T)\)。由\cref{g1:thm:element-in-maximal-torus}，取包含 \(a\) 的极大环面 \(S\)。令 \(K=Z_G(a)\)，并取其单位连通分支 \(K^0\)。它是紧连通 Lie 群：\(K\) 是 \(G\) 的闭子群，而连通分支也闭。

	由于 \(a\) 与 \(T\) 中所有元素交换，\(T\subseteq K\)；由于 \(a\in S\) 且 \(S\) 交换，\(S\subseteq K\)。两环面连通且包含单位元，所以其实都包含于 \(K^0\)。它们在 \(K^0\) 中仍极大，因为任何更大的环面也会是 \(G\) 中更大的环面。将\cref{g1:thm:maximal-tori-conjugate}应用于紧连通群 \(K^0\)，得到 \(k\in K^0\) 使 \(kSk^{-1}=T\)。但 \(k\in Z_G(a)\)，故 \(a=kak^{-1}\in kSk^{-1}=T\)。
\end{proof}

等式 \(\operatorname{Lie}(Z_G(T))=\mathfrak t\) 确定的是中心化子的单位连通分支。证明进一步在 \(Z_G(a)^0\) 中使用共轭定理，排除了环面中心化子的额外分支。对单个元素的中心化子，额外分支仍可能真实存在。

\subsection{环面内共轭与正规化子}
\label{g1:subsec:1903:04}
共轭定理讨论整个子群能否相互移到一起。另一个问题是：已经放在同一个环面里的两个元素，什么时候仍在全群中共轭？环面交换，所以只用环面自身共轭不会改变任何元素，必须考虑那些把环面整体保持住、但不一定逐点固定它的变换。

\begin{definition}\label{g1:def:torus-normalizer}
	设 \(T\) 为 Lie 群 \(G\) 的环面。其\term[zhengguihuazi]{正规化子}[normalizer]定义为
	\[
		N_G(T)\coloneqq\{g\in G:gTg^{-1}=T\}.
	\]
	它通过共轭作用于 \(T\)，通过伴随作用于 \(\mathfrak t\)。中心化子 \(Z_G(T)\) 是这个共轭作用的核。\symidx{\(N_G(T),Z_G(T)\)：环面的正规化子与中心化子}
\end{definition}

\begin{proposition}[环面内共轭的判据]\label{g1:prop:torus-conjugacy-normalizer}
	设 \(G\) 为紧连通实 Lie 群，\(T\subseteq G\) 为极大环面。
	\begin{enumerate}
		\item 对 \(a,b\in T\)，它们在 \(G\) 中共轭，当且仅当存在 \(n\in N_G(T)\) 使 \(nan^{-1}=b\)。
		\item 对 \(X,Y\in\mathfrak t\)，它们在 \(\operatorname{Ad}(G)\) 下属于同一轨道，当且仅当存在 \(n\in N_G(T)\) 使 \(\operatorname{Ad}_nX=Y\)。
	\end{enumerate}
\end{proposition}
\begin{proof}
	1. 若 \(b=gag^{-1}\)，则 \(T\) 与 \(gTg^{-1}\) 都包含于 \(Z_G(b)\)：前者与 \(b\in T\) 交换，后者与 \(b\in gTg^{-1}\) 交换。这两个环面连通，故都在 \(K=Z_G(b)^0\) 中；它们仍是 \(K\) 的极大环面。\(K\) 紧且连通，所以共轭定理给出 \(k\in K\)，使 \(k(gTg^{-1})k^{-1}=T\)。令 \(n=kg\)，便有 \(n\in N_G(T)\)，并且
	\[
		nan^{-1}=kbk^{-1}=b.
	\]
	反向因 \(N_G(T)\subseteq G\) 直接成立。

	2. 若 \(Y=\operatorname{Ad}_gX\)，令
	\[
		K=\{h\in G:\operatorname{Ad}_hY=Y\}^0.
	\]
	括号中的集合是光滑作用的闭稳定子群，所以 \(K\) 是紧连通 Lie 群。由于 \(Y\in\mathfrak t\)，\(T\) 的伴随作用固定 \(Y\)；由于 \(X\in\mathfrak t\)，\(gTg^{-1}\) 的伴随作用也固定 \(Y\)。两环面均连通，所以都包含于 \(K\)，且仍极大。再选 \(k\in K\)，使 \(k(gTg^{-1})k^{-1}=T\)，并令 \(n=kg\)。则 \(n\in N_G(T)\) 且 \(\operatorname{Ad}_nX=\operatorname{Ad}_kY=Y\)。反向同样直接成立。
\end{proof}

因此，每个群元素都能先通过共轭放进固定环面 \(T\)，放进去以后剩余的共轭关系由 \(N_G(T)\) 决定；每个 Lie 代数向量也有对应的陈述。两种分类仍不相同，因为群元素还要经过指数核格的识别。例如在圆周群中，\(0\) 与 \(2\pi\mathrm{i}\) 是不同的 Lie 代数共轭类，却指数到同一个群元素。下一节将把正规化子中逐点固定 \(T\) 的部分除去，得到真正起作用的有限群。

% !TeX root = ../../Book1.tex
% 纲要标题：### 19.4 Weyl 群与根空间
% 纲要位置：几何学计划纲要.md，第 867 行。
% 正文状态：初稿；尚未正式审核。

% 纲要内容开始
% 
% * 正规化子模中心化子
% * 根空间与根的整性
% * 根平面、反射及 Weyl 室
% * 整数格、根格和整体群的区别
% 
% 纲要内容结束

\section{Weyl 群与根空间}
\label{g1:sec:1904}

把元素共轭入一个固定极大环面以后，仍可能有若干环面内的代表。例如，对角酉矩阵可以重新排列对角元；环面本身的共轭作用却是平凡的，不能记录这种排列。本节先把这些剩余对称组成有限群，再研究它为何与一组实超平面的反射有关。始终设 \(G\) 为紧连通实 Lie 群，\(T\) 为极大环面，\(\mathfrak t=\operatorname{Lie}(T)\)，并在 \(\mathfrak g\) 上固定一个伴随不变正定内积 \(B\)。

\subsection{正规化子模中心化子}
\label{g1:subsec:1904:01}
\begin{definition}\label{g1:def:weyl-group}
	设 \(G\) 为紧连通实 Lie 群，\(T\) 为其极大环面。其\term[weylqun]{Weyl 群}[Weyl group]是下列商群。\footnote{Hermann Weyl（1885--1955），德国数学家，在 Lie 群表示论与微分几何方面作出重要贡献。}
	\[
		W=W(G,T)\coloneqq N_G(T)/Z_G(T)=N_G(T)/T.
	\]
	最后一个等号来自\cref{g1:prop:maximal-torus-centralizer}。它通过共轭作用于 \(T\)，通过伴随作用于 \(\mathfrak t\)。\symidx{\(W(G,T)\)：极大环面的 Weyl 群}
\end{definition}

模掉中心化子的理由是它恰好逐点固定环面。若 \(n\) 与 \(nt\) 代表同一个商类，\(t\in T\) 的共轭作用在 \(T\) 上为恒等，故得到的作用不依赖代表。\(T=\exp\mathfrak t\) 又说明：一个元素逐点固定 \(\mathfrak t\)，当且仅当逐点固定 \(T\)。因此这两个作用都忠实。

\begin{proposition}[Weyl 群的有限性]\label{g1:prop:weyl-group-finite}
	设 \(G\) 为紧连通实 Lie 群，\(T\) 为极大环面，\(\mathfrak g=\operatorname{Lie}(G)\)、\(\mathfrak t=\operatorname{Lie}(T)\)。固定 \(\mathfrak g\) 上的 \(\operatorname{Ad}(G)\) 不变正定内积 \(B\)，并在 \(\mathfrak t\) 上取其限制。令 \(W(G,T)=N_G(T)/Z_G(T)\)。则 \(W(G,T)\) 是有限群，在 \(\mathfrak t\) 上忠实地作用为正交变换，并保持整数格 \(\Lambda_T=\ker(\exp_G|_{\mathfrak t})\)。
\end{proposition}
\begin{proof}
	正规化子是闭集：若 \(n_j\to n\) 且 \(n_jTn_j^{-1}=T\)，对每个 \(t\in T\) 取极限得到 \(ntn^{-1}\in T\)；再对 \(n_j^{-1}\) 使用同一论证，便得反向包含。故它是紧 Lie 子群。

	若 \(X\in\operatorname{Lie}(N_G(T))\)，则 \(\operatorname{Ad}_{\exp(sX)}\) 保持 \(\mathfrak t\)，微分得 \([X,H]\in\mathfrak t\) 对所有 \(H\in\mathfrak t\) 成立。另一方面，对于 \(K\in\mathfrak t\)，内积不变性给出
	\[
		B([X,H],K)=B(X,[H,K])=0.
	\]
	所以 \([X,H]\) 同时属于 \(\mathfrak t\) 及其正交补，只能为零。由\cref{g1:prop:maximal-torus-lie-centralizer}，\(X\in\mathfrak t\)。反向包含来自 \(T\subseteq N_G(T)\)，故商 \(N_G(T)/T\) 的维数为零。它既紧又离散，因而有限。

	忠实性已由作用的核说明；正交性来自 \(B\) 的伴随不变性。若 \(H\in\Lambda_T\)，则指数自然性给出
	\[
		\exp_G(\operatorname{Ad}_nH)=n\exp_G(H)n^{-1}=e.
	\]
	故 \(\operatorname{Ad}_n\) 保持 \(\Lambda_T\)，对 \(n^{-1}\) 再用一次得到集合相等。
\end{proof}

\cref{g1:prop:torus-conjugacy-normalizer}现在可以简洁地表述为：\(G\) 的共轭类由 \(T/W\) 参数化，\(\mathfrak g\) 的伴随轨道由 \(\mathfrak t/W\) 参数化。这里先指集合的一一对应；每个轨道与环面相交，而交点恰是同一 \(W\) 轨道。两种参数空间仍有指数核格的差别，不能把 \(T/W\) 与 \(\mathfrak t/W\) 混同。

\subsection{根空间与根的整性}
\label{g1:subsec:1904:02}
\cref{g1:prop:torus-weight-decomposition,g1:prop:maximal-torus-lie-centralizer}已经把互相交换的斜对称算子同时分解为
\begin{equation}\label{g1:eq:root-space-decomposition}
	\mathfrak g_{\mathbb C}=\mathfrak t_{\mathbb C}\oplus
	\bigoplus_{\alpha\in R}\mathfrak g_\alpha,
	\qquad
	[H,Z]=\mathrm{i}\alpha(H)Z\quad(Z\in\mathfrak g_\alpha).
\end{equation}
这里 \(R\subseteq\mathfrak t^*\setminus\{0\}\) 有限。伴随表示的这些非零无穷小权称为\term[gen]{根}[root]，\(\mathfrak g_\alpha\) 称为\term[genkongjian]{根空间}[root space]。我们使用实值的根，因而特征值中显式写出 \(\mathrm{i}\)。取共轭将 \(\mathfrak g_\alpha\) 变成 \(\mathfrak g_{-\alpha}\)；根从一开始就成正负对出现。\symidx{\(R,\mathfrak g_\alpha\)：根集与复根空间}

对于 \(Z\in\mathfrak g_\alpha\)、\(U\in\mathfrak g_\beta\)，Jacobi 恒等式给出
\[
	[H,[Z,U]]=[[H,Z],U]+[Z,[H,U]]
	=\mathrm{i}(\alpha+\beta)(H)[Z,U].
\]
因此 \([\mathfrak g_\alpha,\mathfrak g_\beta]\subseteq\mathfrak g_{\alpha+\beta}\)，其中约定 \(\mathfrak g_0=\mathfrak t_{\mathbb C}\)，不存在的权空间为零。权分解与 Lie 括号相容，正是根比一般表示的权更受约束的原因。下一步要证明每个非零根空间实际上只有一维；一般环面表示的权空间没有这个性质。

\ProseParagraphBreak
反复取括号会沿一个根方向升降特征值。引入相互配合的升、降与权算子，把它们的作用组织成有限链，就能控制根空间的维数，并从链的长度读出整数条件。

\begin{lemma}[升降算子的有限链]\label{g1:lem:sl2-weight-chain}
	设有限维复 Hermitian 空间 \(V\) 上的算子 \(E,F,L\) 满足
	\[
		[L,E]=2E,\qquad [L,F]=-2F,\qquad [E,F]=L,
		\qquad E^*=F,\quad L^*=L.
	\]
	则 \(V\) 是若干相互正交的不变子空间的直和。每个非零子空间有基
	\(v,Fv,\ldots,F^nv\)，其中 \(n\) 是非负整数，且
	\[
		Ev=0,\quad Lv=nv,\quad
		L(F^kv)=(n-2k)F^kv,\quad
		E(F^kv)=k(n-k+1)F^{k-1}v.
	\]
	特别地，\(L\) 的全部特征值都是整数。
\end{lemma}
\begin{proof}
	自伴算子 \(L\) 有正交特征分解。取最大特征值 \(\lambda\) 的非零特征向量 \(v\)。关系 \([L,E]=2E\) 表明 \(Ev\) 若非零就具有特征值 \(\lambda+2\)，所以 \(Ev=0\)。由交换关系归纳得到
	\[
		L(F^kv)=(\lambda-2k)F^kv,\qquad
		E(F^kv)=k(\lambda-k+1)F^{k-1}v.
	\]
	非零的 \(F^kv\) 具有互不相同的特征值，故有限维性保证存在最大整数 \(n\ge0\) 使 \(F^nv\ne0\)，而 \(F^{n+1}v=0\)。在第二式中取 \(k=n+1\)，得到
	\(0=(n+1)(\lambda-n)F^nv\)，故 \(\lambda=n\)。这些向量张成的空间在三个算子下都不变；由于 \(E^*=F\)、\(L^*=L\)，其正交补也不变。例如 \(u\) 在正交补内、\(w\) 在链空间内时，\(\langle Eu,w\rangle=\langle u,Fw\rangle=0\)。在正交补中重复，维数每次严格下降，最终得到所述分解。
\end{proof}

\begin{theorem}[根空间与根的整性]\label{g1:thm:compact-root-spaces}
	设 \(G\) 为紧连通实 Lie 群，\(T\) 为极大环面，\(\mathfrak g=\operatorname{Lie}(G)\)、\(\mathfrak t=\operatorname{Lie}(T)\)。固定 \(\mathfrak g\) 上的 \(\operatorname{Ad}(G)\) 不变正定实内积 \(B\)。采用\eqref{g1:eq:root-space-decomposition}的约定，对 \(\alpha\in\mathfrak t^*\) 令
	\[
		\mathfrak g_\alpha=\{Z\in\mathfrak g_{\mathbb C}:[H,Z]=\mathrm{i}\alpha(H)Z\text{ 对每个 }H\in\mathfrak t\},
		\qquad R=\{\alpha\ne0:\mathfrak g_\alpha\ne0\}.
	\]
	则每个 \(\mathfrak g_\alpha\) 的复维数为一。以 \(B\) 识别 \(\mathfrak t\) 与其对偶，并以 \((\alpha,\beta)\) 记诱导的对偶内积，则
	\[
		\frac{2(\beta,\alpha)}{(\alpha,\alpha)}\in\mathbb Z
		\qquad(\alpha,\beta\in R).
	\]
	若两根 \(\alpha,\beta\) 共线，则 \(\beta=\alpha\) 或 \(\beta=-\alpha\)。
\end{theorem}
\begin{proof}
	把 \(B\) 复双线性延拓为 \(B_{\mathbb C}\)，令 Hermitian 内积为 \(\langle Z,U\rangle=B_{\mathbb C}(Z,\overline U)\)。固定根 \(\alpha\)，取范数一的 \(Z\in\mathfrak g_\alpha\)，并记 \(A_\alpha\in\mathfrak t\) 为满足 \(B(A_\alpha,H)=\alpha(H)\) 的向量，\(a^2=(\alpha,\alpha)=B(A_\alpha,A_\alpha)>0\)。由括号的权相加性质，\([Z,\overline Z]\in\mathfrak t_{\mathbb C}\)；对 \(H\in\mathfrak t\)，计算
	\[
		B_{\mathbb C}([Z,\overline Z],H)
		=B_{\mathbb C}(Z,[\overline Z,H])
		=\mathrm{i}\alpha(H).
	\]
	所以 \([Z,\overline Z]=\mathrm{i}A_\alpha\)。令
	\begin{equation}\label{g1:eq:root-sl2-triple}
		e=\frac{\sqrt2}{a}Z,\qquad
		f=-\frac{\sqrt2}{a}\overline Z,\qquad
		h=-\frac{2\mathrm{i}}{a^2}A_\alpha.
	\end{equation}
	直接代入得 \([e,f]=h\)、\([h,e]=2e\)、\([h,f]=-2f\)。又有 \(\overline e=-f\)、\(\overline h=-h\)。内积不变性说明 \((\operatorname{ad}_u)^*=-\operatorname{ad}_{\overline u}\)，故 \(E=\operatorname{ad}_e\)、\(F=\operatorname{ad}_f\)、\(L=\operatorname{ad}_h\) 满足前一引理。在 \(\mathfrak g_\beta\) 上，\(L\) 的特征值是
	\[
		\frac{2\beta(A_\alpha)}{a^2}
		=\frac{2(\beta,\alpha)}{(\alpha,\alpha)},
	\]
	所以它是整数。

	为控制 \(\mathfrak g_\alpha\) 的维数，只看不变子空间
	\[
		V_\alpha=\mathfrak t_{\mathbb C}\oplus
		\bigoplus_{\substack{k\in\mathbb Z\setminus\{0\}\\k\alpha\in R}}
		\mathfrak g_{k\alpha}.
	\]
	它的 \(L\) 特征值都是偶数，零特征空间恰为 \(\mathfrak t_{\mathbb C}\)。升算子在这个零特征空间上满足 \(E(H)=-\mathrm{i}\alpha(H)e\)，所以其秩恰为一。引理中的每条非平凡偶数链都经过权零与权二，且从权零到权二的映射非零；因此这种链只能有一条。另一方面，\(\operatorname{span}_{\mathbb C}\{e,h,f\}\) 已是一条最高权为二的链，其正交补不变。于是余下的链全为平凡链。由此得到 \(\dim_{\mathbb C}\mathfrak g_\alpha=1\)，并且 \(k\alpha\notin R\) 对所有整数 \(|k|\ge2\) 成立。

	最后若 \(\beta=c\alpha\)、\(c>0\)，上面的整性分别用于 \((\beta,\alpha)\) 与 \((\alpha,\beta)\)，得到 \(2c\) 和 \(2/c\) 都是正整数；它们的乘积为四，故 \(c\in\{1/2,1,2\}\)。\(c=2\) 已被排除，\(c=1/2\) 则对根 \(\beta\) 使用同一结论排除。负倍数的情形先将 \(\beta\) 换为 \(-\beta\)，所以只剩 \(\pm\alpha\)。
\end{proof}

这里还可以把根与群元素的正则性联系起来。对 \(g\in G\)，若 \(\dim Z_G(g)=\operatorname{rank}G\)，则称该元素\term[qunyuansuzhengzexing]{正则}[regular]。对 \(t\in T\)，权分解给出
\[
	\operatorname{Lie}(Z_G(t))_{\mathbb C}
	=\mathfrak t_{\mathbb C}\oplus
	\bigoplus_{\substack{\alpha\in R\\\chi_\alpha(t)=1}}\mathfrak g_\alpha,
	\qquad \chi_\alpha(\exp H)=e^{\mathrm{i}\alpha(H)}.
\]
因此正则性等价于每个 \(\chi_\alpha(t)\ne1\)。这只控制中心化子的维数，即其单位分支；中心化子仍可能有其他分支。与 Lie 代数条件 \(\alpha(H)\ne0\) 相比，群元素 \(\exp H\) 要避开的是全部 \(\alpha(H)\in2\pi\mathbb Z\)。

\subsection{根平面、反射及 Weyl 室}
\label{g1:subsec:1904:03}
复根空间为一维，但原来的 Lie 代数是实空间。每一对 \(\{\alpha,-\alpha\}\) 恢复一个实二维平面
\[
	P_\alpha=(\mathfrak g_\alpha\oplus\mathfrak g_{-\alpha})\cap\mathfrak g,
	\qquad P_\alpha=P_{-\alpha},
\]
称为相应的\term[genpingmian]{根平面}[root plane]。取上一证明中范数一的 \(Z\in\mathfrak g_\alpha\)，则
\[
	u=\frac{Z+\overline Z}{\sqrt2},\qquad
	v=\frac{\mathrm{i}(Z-\overline Z)}{\sqrt2}
\]
构成 \(P_\alpha\) 的实正交单位基。这里 \(B_{\mathbb C}(Z,Z)=0\)：对取值 \(\alpha(H)\ne0\) 的 \(H\) 应用内积不变性，得到 \(2\mathrm{i}\alpha(H)B_{\mathbb C}(Z,Z)=0\)，再与 \(B_{\mathbb C}(Z,\overline Z)=1\) 合用便可核对正交与单位长度。括号为
\[
	[H,u]=\alpha(H)v,\qquad [H,v]=-\alpha(H)u.
\]
所以 \(\operatorname{Ad}_{\exp H}\) 在这个平面上旋转角度 \(\alpha(H)\)。选取每一对根中的一个组成集合 \(R_+\)，就有实正交分解
\[
	\mathfrak g=\mathfrak t\oplus\bigoplus_{\alpha\in R_+}P_\alpha,
	\qquad \dim_{\mathbb R}G=\operatorname{rank}G+2|R_+|.
\]
根平面在 \(\mathfrak g\) 内，而方程 \(\alpha(H)=0\) 的超平面在 \(\mathfrak t\) 内；它们是两个不同空间里的对象。前者表示被旋转的方向，后者表示旋转的无穷小速度为零的参数。

\ProseParagraphBreak
令
\begin{equation}\label{g1:eq:coroot-reflection}
	\alpha^\vee\coloneqq\frac{2A_\alpha}{(\alpha,\alpha)},
	\qquad
	s_\alpha(H)\coloneqq H-\alpha(H)\alpha^\vee.
\end{equation}
向量 \(\alpha^\vee\) 称为\term[yugen]{余根}[coroot]。\symidx{\(\alpha^\vee,s_\alpha\)：余根与根反射}由于 \(\alpha(\alpha^\vee)=2\)，\(s_\alpha\) 固定 \(\ker\alpha\)，并把它的正交方向 \(\alpha^\vee\) 变号，所以它就是关于 \(\ker\alpha\) 的正交反射。定义一个线性反射尚不足以证明它来自群的共轭作用；下面给出实现它的具体群元素。

\begin{proposition}[根反射的群内实现]\label{g1:prop:root-reflection-realized}
	设 \(G\) 为紧连通实 Lie 群，\(T\) 为极大环面，\(\mathfrak g=\operatorname{Lie}(G)\)、\(\mathfrak t=\operatorname{Lie}(T)\)。固定 \(\mathfrak g\) 上的 \(\operatorname{Ad}(G)\) 不变正定内积 \(B\)，以 \((\alpha,\beta)\) 记 \(B|_{\mathfrak t}\) 诱导的对偶内积。令 \(R\) 为所有满足下列条件的非零实泛函 \(\alpha\in\mathfrak t^*\) 的集合：存在 \(0\ne Z\in\mathfrak g_{\mathbb C}\)，使 \([H,Z]=\mathrm{i}\alpha(H)Z\) 对每个 \(H\in\mathfrak t\) 成立。对 \(\alpha\in R\)，以 \(B(H_\alpha,H)=\alpha(H)\)（\(H\in\mathfrak t\)）定义 \(H_\alpha\)，并令 \(\alpha^\vee=2H_\alpha/B(H_\alpha,H_\alpha)\)。令 \(s_\alpha(H)=H-\alpha(H)\alpha^\vee\)。则每个根 \(\alpha\) 的反射 \(s_\alpha\) 都属于 \(W(G,T)\) 在 \(\mathfrak t\) 上的作用。整个 \(W(G,T)\) 保持根集；其对根的作用为 \((w\alpha)(H)=\alpha(w^{-1}H)\)。特别地，
	\[
		s_\alpha\beta
		=\beta-\frac{2(\beta,\alpha)}{(\alpha,\alpha)}\alpha
		\in R\qquad(\alpha,\beta\in R).
	\]
\end{proposition}
\begin{proof}
	沿用\eqref{g1:eq:root-sl2-triple}中的 \(e,f,h\)，令
	\[
		x=e-f,\qquad y=\mathrm{i}(e+f),\qquad k=\mathrm{i}h=\alpha^\vee.
	\]
	它们在复共轭下不变，故属于实 Lie 代数，且括号满足
	\[
		[x,y]=2k,\qquad [y,k]=2x,\qquad[k,x]=2y.
	\]
	因此 \(\operatorname{ad}_x\) 在 \(\operatorname{span}_{\mathbb R}\{y,k\}\) 上的矩阵为
	\(\begin{psmallmatrix}0&-2\\2&0\end{psmallmatrix}\)。由\cref{g1:prop:exponential-naturality}，
	\[
		n_\alpha=\exp_G(\tfrac\pi2x)
		\quad\Longrightarrow\quad
		\operatorname{Ad}_{n_\alpha}k=-k.
	\]
	若 \(H\in\ker\alpha\)，则它与 \(e,f\) 都交换，因而与 \(x\) 交换，得到 \(\operatorname{Ad}_{n_\alpha}H=H\)。将一般 \(H\in\mathfrak t\) 写成
	\[
		H=\left(H-\frac{\alpha(H)}2k\right)+\frac{\alpha(H)}2k,
	\]
	括号内属于 \(\ker\alpha\)，从而 \(\operatorname{Ad}_{n_\alpha}H=H-\alpha(H)k=s_\alpha(H)\)。该作用保持 \(\mathfrak t\)，故指数自然性和 \(T=\exp\mathfrak t\) 给出 \(n_\alpha Tn_\alpha^{-1}=T\)。所以它确实代表一个 Weyl 群元素。

	对任意 \(n\in N_G(T)\)，令 \(w=\operatorname{Ad}_n|_{\mathfrak t}\)。若 \(Z\in\mathfrak g_\beta\)，则
	\[
		[H,\operatorname{Ad}_nZ]
		=\operatorname{Ad}_n[w^{-1}H,Z]
		=\mathrm{i}\beta(w^{-1}H)\operatorname{Ad}_nZ.
	\]
	因 \(\operatorname{Ad}_nZ\ne0\)，\(w\beta\) 仍是根。取 \(w=s_\alpha=w^{-1}\)，代入\eqref{g1:eq:coroot-reflection}，便得最后的公式。
\end{proof}

至此，根集已具有有限性、成对性、整性和反射不变性，且一条直线上只有一对根。在 \(\operatorname{span}_{\mathbb R}R\) 中，这些性质正组成一个约化的晶体根系。根也未必张成全部 \(\mathfrak t^*\)：其公共零空间恰为 \(\mathfrak g\) 的中心。事实上，\(H\in\mathfrak t\) 与每个根空间交换当且仅当全部 \(\alpha(H)=0\)，而零权空间本来就与 \(H\) 交换；反过来，中心元素先与 \(\mathfrak t\) 交换，故属于 \(\mathfrak t\)。

\begin{definition}\label{g1:def:weyl-chamber}
	设 \(T\) 为紧连通 Lie 群的极大环面，\(R\) 为其根集。开集
	\[
		\mathfrak t_{\mathrm{reg}}
		=\mathfrak t\setminus\bigcup_{\alpha\in R}\ker\alpha
	\]
	的各连通分支称为\term[weylshi]{Weyl 室}[Weyl chamber]。室的闭包中余维一的边界部分称为它的墙。
\end{definition}

选定 \(H_0\in\mathfrak t_{\mathrm{reg}}\)，令 \(R_+=\{\alpha:\alpha(H_0)>0\}\)。包含 \(H_0\) 的室就是
\[
	C=\{H:\alpha(H)>0\;\text{对全部 }\alpha\in R_+\}.
\]
一方面，一个连通分支内每个根的符号不能变化；另一方面，右侧是有限个开半空间的交，凸且连通，所以恰为这个分支。室是开凸多面锥；若存在中心方向，它沿这些方向不受限制。若 \(G\) 本身是环面，根集为空，唯一的室就是整个 \(\mathfrak t\)。

\begin{proposition}[相邻室与反射]\label{g1:prop:chamber-reflection-transitivity}
	设 \(G\) 为紧连通实 Lie 群，\(T\) 为极大环面。\(W(G,T)\) 置换 Weyl 室。由所有根反射生成的子群在这些室上作用传递；隔着同一墙相邻的两个室由关于该墙的反射交换。
\end{proposition}
\begin{proof}
	Weyl 群保持根集，所以保持根超平面的并，进而置换其补集的连通分支。设两个室沿 \(\ker\alpha\) 的一片墙相邻。在这片墙上取不属于其他不同根超平面的点 \(q\)，再取充分小的球，使球中只经过这一个根超平面。反射 \(s_\alpha\) 固定 \(q\)，交换球内的两个半球，又保持整个超平面排列；因此它把这两个半球所在的室相互交换。

	任取两个室 \(C,D\)，选 \(p\in C\)，再在 \(D\) 中选 \(q\)，使线段 \([p,q]\) 不经过两个不同根超平面的交。这样的 \(q\) 存在：对每个余维至少为二的交空间 \(L\)，所有可能使线段碰到 \(L\) 的端点都落在包含 \(p\) 与 \(L\) 的一个真仿射子空间内；有限个真仿射子空间不能覆盖开集 \(D\)。这条线段只在有限个互不相同的时刻逐个穿过根超平面，依次经过一列相邻室。每一次相邻转移由刚证明的一个根反射实现，将这些反射复合，就把 \(C\) 送到 \(D\)。根集为空时只有一个室，结论同样成立。
\end{proof}

% !TeX root = ../Book1.tex
\begin{figure}[htbp]
	\centering
	\begin{tikzpicture}[font=\small,scale=0.95]
		\begin{scope}
			\fill[GeometryBlue!10] (0,0)--(30:2.25) arc[start angle=30,end angle=90,radius=2.25]--cycle;
			\foreach \a in {30,90,150}{
				\draw[GeometryGray,dashed] (\a:2.4)--(\a:-2.4);
			}
			\foreach \a in {0,60,120,180,240,300}{
				\draw[->,very thick,GeometryBlue] (0,0)--(\a:1.7);
			}
			\node[right] at (0:1.75) {\(\alpha\)};
			\node[above left] at (120:1.75) {\(\beta\)};
			\node[above right] at (60:1.75) {\(\alpha+\beta\)};
			\node at (60:0.8) {\(C\)};
			\node at (0,-2.85) {\(A_2\)：六根，六个室};
		\end{scope}
		\begin{scope}[xshift=6cm]
			\fill[GeometryBlue!10] (0,0)--(0:2.25) arc[start angle=0,end angle=45,radius=2.25]--cycle;
			\foreach \a in {0,45,90,135}{
				\draw[GeometryGray,dashed] (\a:2.4)--(\a:-2.4);
			}
			\foreach \a in {0,90,180,270}{
				\draw[->,very thick,GeometryBlue] (0,0)--(\a:1.35);
			}
			\foreach \x/\y in {1.35/1.35,-1.35/1.35,-1.35/-1.35,1.35/-1.35}{
				\draw[->,very thick,GeometryBlue] (0,0)--(\x,\y);
			}
			\node[below] at (1.45,0) {\(\varepsilon_1\)};
			\node[above left] at (0,1.4) {\(\varepsilon_2\)};
			\node[right] at (1.35,1.35) {\(\varepsilon_1+\varepsilon_2\)};
			\node at (22.5:0.9) {\(C\)};
			\node at (0,-2.85) {\(B_2\)：八根，八个室};
		\end{scope}
	\end{tikzpicture}
	\caption[秩二根系的根与反射墙]{实根借内积画为蓝色向量，灰色虚线为根的核，浅蓝色区域各标出一个 Weyl 室。左图的根等长；右图的长根长度为短根的 \(\sqrt2\) 倍。图示的是 \(\mathfrak t\) 与对偶空间经内积识别后的超平面排列，不是 \(\mathfrak g\) 中的实根平面。}
	\label{g1:fig:root-arrangements}
\end{figure}


\cref{g1:fig:root-arrangements}画出秩二根系的两个模型。图中箭头把根借内积画成向量，虚线才是它们的核；位于根方向的箭头与同根对应的反射墙互相垂直。一般群的根反射与室之间的关系由上面的命题保证；下一节还会从矩阵直接算出经典群的整个 Weyl 群。

\subsection{整数格、根格和整体群的区别}
\label{g1:subsec:1904:04}
根来自伴随表示的微分，主要描述 Lie 代数。要恢复环面中的周期，还必须保留指数核
\(\Lambda_T\subseteq\mathfrak t\)。为了把这些不同层次的数据放在一起，先把环面的全部特征标写成线性泛函。

\begin{proposition}[环面的特征标格]\label{g1:prop:torus-character-lattice}
	设 \(T=\mathfrak t/\Lambda_T\) 为有限维紧连通交换 Lie 群。光滑群同态 \(\chi:T\to S^1\) 与集合
	\[
		X^*(T)\coloneqq\{\lambda\in\mathfrak t^*:
		\lambda(\Lambda_T)\subseteq2\pi\mathbb Z\}
	\]
	中的泛函一一对应，对应公式为 \(\chi_\lambda(\exp H)=e^{\mathrm{i}\lambda(H)}\)。该集合在加法下是 \(\mathfrak t^*\) 中的满秩格。\symidx{\(\Lambda_T,X^*(T)\)：环面的指数核格与特征标格}
\end{proposition}
\begin{proof}
	同态 \(\chi\) 的单位元微分是实线性映射 \(\mathfrak t\to\mathrm{i}\mathbb R\)，可唯一写成 \(H\mapsto\mathrm{i}\lambda(H)\)。指数自然性给出所述公式。对于 \(H\in\Lambda_T\)，必须有 \(e^{\mathrm{i}\lambda(H)}=1\)，这就是整性条件。反过来，满足该条件时，公式在 \(H\) 换成 \(H+L\)、\(L\in\Lambda_T\) 后不变，所以能下降到商环面，且是光滑群同态。指数的满射性保证对应的唯一性。

	取 \(\Lambda_T\) 的格基 \(L_1,\ldots,L_r\)，令 \(\lambda_j(L_k)=2\pi\delta_{jk}\)。上述条件恰好要求 \(\lambda\) 是 \(\lambda_1,\ldots,\lambda_r\) 的整数线性组合，故它们组成 \(X^*(T)\) 的格基。
\end{proof}

这里的特征标是一维酉表示；对于一般有限维表示，常把表示矩阵的迹也称为特征标，一维情形两者恰好一致。上一节已经得到每个根都满足整性，故有\term[genge]{根格}[root lattice]
\[
	Q\coloneqq\sum_{\alpha\in R}\mathbb Z\alpha
	\subseteq X^*(T).
\]
\(Q\) 位于对偶空间，而 \(\Lambda_T\) 位于原空间，不能直接比较包含关系。两者通过取值配对联系：\(\lambda(L)/(2\pi)\in\mathbb Z\)。根格在 \(\operatorname{span}_{\mathbb R}R\) 中是满秩格，因为它是满秩格 \(X^*(T)\) 的离散子群，并张成这个子空间；若有连续中心方向，根格的秩可能小于 \(\dim T\)。

\ProseParagraphBreak
余根属于原空间，因而可以与周期格直接比较；正确的系数是 \(2\pi\)。

\begin{proposition}[余根给出的周期]\label{g1:prop:coroot-period}
	设 \(G\) 为紧连通实 Lie 群，\(T\) 为极大环面，\(\mathfrak g=\operatorname{Lie}(G)\)、\(\mathfrak t=\operatorname{Lie}(T)\)。固定 \(\mathfrak g\) 上的 \(\operatorname{Ad}(G)\) 不变正定内积 \(B\)，以 \((\alpha,\beta)\) 记 \(B|_{\mathfrak t}\) 诱导的对偶内积。令 \(R\) 为所有满足下列条件的非零实泛函 \(\alpha\in\mathfrak t^*\) 的集合：存在 \(0\ne Z\in\mathfrak g_{\mathbb C}\)，使 \([H,Z]=\mathrm{i}\alpha(H)Z\) 对每个 \(H\in\mathfrak t\) 成立。对 \(\alpha\in R\)，以 \(B(H_\alpha,H)=\alpha(H)\)（\(H\in\mathfrak t\)）定义 \(H_\alpha\)，并令 \(\alpha^\vee=2H_\alpha/B(H_\alpha,H_\alpha)\)。根与余根采用\eqref{g1:eq:root-space-decomposition}及\eqref{g1:eq:coroot-reflection}的归一化，并记 \(\Lambda_T=\ker(\exp_G|_{\mathfrak t})\)。则
	\[
		\exp_G(2\pi\alpha^\vee)=e\qquad(\alpha\in R).
	\]
	因此若 \(Q^\vee=\sum_{\alpha\in R}\mathbb Z\alpha^\vee\)，则 \(2\pi Q^\vee\subseteq\Lambda_T\)。
\end{proposition}
\begin{proof}
	在\cref{g1:prop:root-reflection-realized}的证明中，实向量 \(x,y,k=\alpha^\vee\) 满足与下列矩阵相同的括号关系：
	\[
		x_0=\begin{pmatrix}0&1\\-1&0\end{pmatrix},\qquad
		y_0=\begin{pmatrix}0&\mathrm{i}\\\mathrm{i}&0\end{pmatrix},\qquad
		k_0=\begin{pmatrix}\mathrm{i}&0\\0&-\mathrm{i}\end{pmatrix}.
	\]
	它们是 \(\mathfrak{su}(2)\) 的实基，所以 \(x_0\mapsto x\)、\(y_0\mapsto y\)、\(k_0\mapsto k\) 定义了 Lie 代数同态 \(\varphi:\mathfrak{su}(2)\to\mathfrak g\)。这里需要一个整体同态来保持完整的周期，而不仅是局部指数公式。

	每个 \(\operatorname{SU}(2)\) 矩阵唯一写成
	\[
		\begin{pmatrix}z&w\\-\overline w&\overline z\end{pmatrix},
		\qquad |z|^2+|w|^2=1.
	\]
	这由两行酉正交及行列式等于一直接得到；取出 \(z,w\) 给出光滑逆映射，所以 \(\operatorname{SU}(2)\) 微分同胚于 \(S^3\)。由\cref{g1:prop:sphere-simply-connected}，它连通且单连通。\cref{g1:thm:lie-homomorphism-integration}遂把 \(\varphi\) 积分为 \(\Phi:\operatorname{SU}(2)\to G\)。再用指数自然性和对角矩阵指数，得到
	\[
		\exp_G(2\pi\alpha^\vee)
		=\Phi\bigl(\exp_{\operatorname{SU}(2)}(2\pi k_0)\bigr)
		=\Phi(I)=e.
	\]
	余根都在交换子代数 \(\mathfrak t\) 内，其整数线性组合的指数等于相应指数的乘积，故得到所列格包含。
\end{proof}

若根张成 \(\mathfrak t^*\)，还可定义\term[quange]{权格}[weight lattice]及\term[yuquange]{余权格}[coweight lattice]
\[
	P=\{\lambda\in\mathfrak t^*:\lambda(\alpha^\vee)\in\mathbb Z
	\;\text{对所有 }\alpha\in R\},\qquad
	P^\vee=\{H\in\mathfrak t:\alpha(H)\in\mathbb Z
	\;\text{对所有 }\alpha\in R\}.
\]
根张成假设保证这两者都是满秩格。的确，\(Q\) 已是满秩格，而 \(P^\vee\) 是它的整数对偶格；\(Q^\vee\subseteq\Lambda_T/(2\pi)\) 是离散且张成 \(\mathfrak t\) 的子群，故也是满秩格，\(P\) 是其整数对偶格。已有结论合成
\begin{equation}\label{g1:eq:root-lattice-inclusions}
	Q\subseteq X^*(T)\subseteq P,
	\qquad
	Q^\vee\subseteq\frac1{2\pi}\Lambda_T\subseteq P^\vee.
\end{equation}
\symidx{\(Q,P,Q^\vee,P^\vee\)：根格、权格、余根格与余权格}前一组右侧的包含是把 \(\lambda(\Lambda_T)\subseteq2\pi\mathbb Z\) 用于 \(2\pi\alpha^\vee\)；后一组右侧的包含则是把根的整性用于任意周期。若有中心方向，\(Q\)、\(Q^\vee\) 只张成相应的非中心子空间，不能把这些定义不加条件地称为整个空间的满秩格。

\ProseParagraphBreak
根还直接检测群中心，但需要在环面上取指数以后的值。

\begin{proposition}[由根检测中心]\label{g1:prop:center-from-roots}
	设 \(G\) 为紧连通实 Lie 群，\(T\) 为极大环面，\(\mathfrak g=\operatorname{Lie}(G)\)、\(\mathfrak t=\operatorname{Lie}(T)\)。令 \(R\) 为所有满足下列条件的非零实泛函 \(\alpha\in\mathfrak t^*\) 的集合：存在 \(0\ne Z\in\mathfrak g_{\mathbb C}\)，使 \([H,Z]=\mathrm{i}\alpha(H)Z\) 对每个 \(H\in\mathfrak t\) 成立。对 \(\alpha\in R\)，令根特征标 \(\chi_\alpha:T\to S^1\) 满足 \(\chi_\alpha(\exp_G H)=e^{\mathrm{i}\alpha(H)}\)，并记
	\[
		\Lambda_T=\ker(\exp_G|_{\mathfrak t}),\qquad
		P^\vee=\{H\in\mathfrak t:\alpha(H)\in\mathbb Z\text{ 对所有 }\alpha\in R\}.
	\]
	则
	\[
		Z(G)=\{t\in T:\chi_\alpha(t)=1\;\text{对所有 }\alpha\in R\}.
	\]
	若根张成 \(\mathfrak t^*\)，则还有自然群同构
	\[
		Z(G)\cong (2\pi P^\vee)/\Lambda_T.
	\]
\end{proposition}
\begin{proof}
	中心元素与 \(T\) 交换，故由 \(Z_G(T)=T\) 必在 \(T\) 中。对于 \(t\in T\)，其伴随作用在 \(\mathfrak t_{\mathbb C}\) 上恒等，在 \(\mathfrak g_\alpha\) 上乘以 \(\chi_\alpha(t)\)；所以所列根条件等价于 \(\operatorname{Ad}_t\) 在整个 \(\mathfrak g\) 上恒等。指数自然性于是说明 \(t\) 与所有指数交换，而连通群由指数生成，故 \(t\in Z(G)\)。这也证明反向包含。

	写 \(t=\exp H\)，根条件为 \(\alpha(H)\in2\pi\mathbb Z\)，即 \(H\in2\pi P^\vee\)。指数在此子群上的核仍为 \(\Lambda_T\)，所以得到最后的商群描述。
\end{proof}

例如，将 \(G\) 除以有限中心子群 \(C\)，不会改变 Lie 代数、根空间和根的配对，却会把周期格变成
\[
	\Lambda_{T/C}=\{H\in\mathfrak t:\exp_GH\in C\}\supseteq\Lambda_T.
\]
这里 \(C\subseteq T\) 已由中心判据保证。一个原来的特征标能下降到 \(T/C\)，当且仅当它在 \(C\) 上恒为一，故商群的特征标格反而缩小。下一节的 \(\operatorname{SU}(2)\) 与 \(\operatorname{SO}(3)\) 将把这一现象写成 \(2\pi\mathbb Z\) 与 \(\pi\mathbb Z\) 的具体区别。根记录括号和无穷小共轭，周期格记录哪些一参数运动已经闭合；分清这两类数据，才能讨论整体群的中心和覆盖。
% 纲要条目排为编号小节标题，并进入目录与书签；说明段仍保留于纲要注释中。

% !TeX root = ../../Book1.tex
% 纲要标题：### 19.5 经典紧群的完整算例
% 纲要位置：几何学计划纲要.md，第 874 行。
% 正文状态：初稿；尚未正式审核。

% 纲要内容开始
% 
% * 酉群和特殊酉群的环面与根
% * 正交群与辛群的根模型
% * 低秩同构、中心及基本群
% 
% ---
% 
% 纲要内容结束

\section{经典紧群的完整算例}
\label{g1:sec:1905}

抽象理论把共轭问题集中到一个环面，再把环面内剩余的识别交给 Weyl 群。对矩阵群，这两步分别对应选择正交或酉基，以及重新排列特征方向。本节把环面、根、整数格、中心和基本群放在同一组模型中计算。根描述伴随作用的无穷小旋转，整数格则记录哪些旋转已经回到单位元；两者同时出现，才不会把具有相同 Lie 代数的群误认为同一个群。

\subsection{酉群和特殊酉群的环面与根}
\label{g1:subsec:1905:01}
先看最熟悉的酉对角化。记 \(\mathfrak u(n)=\{X:X^*=-X\}\)，并令 \(E_{jk}\) 为第 \((j,k)\) 项等于一、其余项为零的矩阵。

\begin{proposition}[酉群的环面与根]\label{g1:prop:unitary-torus-roots}
	设 \(n\ge2\)。在 \(\operatorname{U}(n)\) 中，
	\[
		T=\{\operatorname{diag}(e^{\mathrm{i}\theta_1},\ldots,e^{\mathrm{i}\theta_n}):\theta_j\in\mathbb R\}
	\]
	是极大环面，其 Lie 代数用 \(H(\theta)=\mathrm{i}\operatorname{diag}(\theta_1,\ldots,\theta_n)\) 识别为 \(\mathbb R^n\)。若 \(\varepsilon_j(H(\theta))=\theta_j\)，则
	\[
		\Lambda=2\pi\mathbb Z^n,\qquad
		R=\{\varepsilon_j-\varepsilon_k:j\ne k\},\qquad
		W\cong S_n.
	\]
	在 \(\operatorname{SU}(n)\) 中，相应极大环面为 \(T_0=T\cap\operatorname{SU}(n)\)，Lie 代数及整数格分别是
	\[
		\mathfrak t_0=\{H(\theta):\textstyle\sum_j\theta_j=0\},\qquad
		\Lambda_0=\{2\pi(m_1,\ldots,m_n):m_j\in\mathbb Z,\ \textstyle\sum_jm_j=0\}.
	\]
	其根是上述根在 \(\mathfrak t_0\) 上的限制，Weyl 群仍为 \(S_n\)。
\end{proposition}
\begin{proof}
	酉矩阵有酉特征基。这里所需的谱论可以从一个单位特征向量开始归纳：复特征值存在，酉性使其模为一；该向量的正交补在矩阵及其逆下都不变，因而可继续选基。于是每个酉矩阵酉共轭于 \(T\)。改变对角元的角度给出从单位矩阵到它的路径，再作同一共轭，说明 \(\operatorname{U}(n)\) 连通。对于特殊酉矩阵，特征角之和属于 \(2\pi\mathbb Z\)；给其中一个角减去这个整数倍后，可使角度和恰为零，路径 \(\operatorname{diag}(e^{\mathrm{i}t\theta_j})\) 就始终在 \(\operatorname{SU}(n)\) 中。对角化所用的酉矩阵也可右乘一个对角酉矩阵，将其行列式调成一而不改变对角化。因此 \(\operatorname{SU}(n)\) 连通，且每个元素都在该群内共轭入 \(T_0\)。两群由矩阵方程定义为有界闭集，所以紧。

	\(T_0\) 可由前 \(n-1\) 个相位自由参数化，最后一个取其乘积的逆，故是连通环面。在两种情形，与所有这些对角矩阵交换的矩阵必须是对角矩阵：对第 \((j,k)\) 项选取第 \(j,k\) 个相位不同的环面元即可。于是两群中的中心化子分别就是 \(T,T_0\)，任何包含它们的环面也不能更大。这直接验证了极大性。

	\(\mathfrak u(n)\) 的复化是 \(\mathfrak{gl}(n,\mathbb C)\)，而 \(\mathfrak{su}(n)\) 的复化是 \(\mathfrak{sl}(n,\mathbb C)\)。逐项计算得
	\[
		[H(\theta),E_{jk}]=\mathrm{i}(\theta_j-\theta_k)E_{jk}.
	\]
	故每个非对角矩阵单位给出所列根，零权空间是对角子代数；取迹零只删去中心的一维零权方向。指数是逐项取指数，所以整数格正是所写的集合。

	正规化环面的酉矩阵必须置换其共同特征直线 \(\mathbb Ce_j\)，因而是一个置换矩阵与对角酉矩阵的乘积；反之这样的矩阵确实正规化环面。模掉对角部分只剩置换。在特殊酉群中，任意置换矩阵都能乘以对角酉矩阵将行列式调成一，仍实现相同的置换。因此两个商群都是 \(S_n\)。
\end{proof}

上述根系在迹零超平面内通常记作 \(A_{n-1}\)。一个 Weyl 室可取 \(\theta_1>\theta_2>\cdots>\theta_n\)，在特殊酉情形还要求和为零。相邻两角相等的墙记录两个特征方向不能再由伴随作用区分。\(\operatorname{U}(n)\) 还保留所有角同时增加的中心方向，所以它的根不张成整个 \(\mathfrak t^*\)。

\begin{example}\label{g1:exm:su2-so3-lattices}
	在 \(\operatorname{SU}(2)\) 中写 \(H(t)=\operatorname{diag}(\mathrm{i}t,-\mathrm{i}t)\)。正根为 \(\alpha(H(t))=2t\)，而 \(\exp H(t)=I\) 当且仅当 \(t\in2\pi\mathbb Z\)。商群 \(\operatorname{SU}(2)/\{\pm I\}\) 使用同一个 Lie 代数与同一条根，却在 \(t\in\pi\mathbb Z\) 时就回到单位元。因此根空间完全相同，整数格却扩大一倍。前一个环面的特征标为 \(e^{\mathrm{i}mt}\)，\(m\in\mathbb Z\)；能下降到商环面的特征标还必须在 \(t=\pi\) 取值一，即 \(m\) 必须为偶数。这个商群与 \(\operatorname{SO}(3)\) 的同构将在本节后面直接构造。
\end{example}

\subsection{正交群与辛群的根模型}
\label{g1:subsec:1905:02}
正交群的实特征方向未必是一维。非实特征值成共轭对，对应实二维旋转平面；这使它的环面参数成对出现。令
\[
	R(\theta)=\begin{pmatrix}\cos\theta&-\sin\theta\\\sin\theta&\cos\theta\end{pmatrix},\qquad
	J=\begin{pmatrix}0&-1\\1&0\end{pmatrix}.
\]

\begin{proposition}[旋转群的根模型]\label{g1:prop:orthogonal-torus-roots}
	令 \(J=\begin{pmatrix}0&-1\\1&0\end{pmatrix}\)。在 \(\operatorname{SO}(2r)\)（\(r\ge2\)）中取由旋转块 \(\exp(\theta_jJ)\) 组成的分块对角环面，在 \(\operatorname{SO}(2r+1)\)（\(r\ge1\)）中另加一个末尾的 \(1\) 块。以 \(H(\theta)=\operatorname{diag}(\theta_1J,\ldots,\theta_rJ)\) 表示其 Lie 代数，奇数维时另加零块，并置 \(\varepsilon_j(H(\theta))=\theta_j\)。这两个环面都是极大的，整数格均为 \(2\pi\mathbb Z^r\)，其根与 Weyl 群为
	\[
		\begin{array}{c|c|c}
		G&R&W\\ \hline
		\operatorname{SO}(2r+1)&\{\pm\varepsilon_j\pm\varepsilon_k:j<k\}\cup\{\pm\varepsilon_j\}&(\mathbb Z/2)^r\rtimes S_r\\
		\operatorname{SO}(2r)&\{\pm\varepsilon_j\pm\varepsilon_k:j<k\}&(\mathbb Z/2)^{r-1}\rtimes S_r.
		\end{array}
	\]
	偶数维一行中的符号变换要求改变偶数个坐标的符号。
\end{proposition}
\begin{proof}
	\cref{g1:thm:orthogonal-blocks}的实正交标准分块给出旋转块及 \(1,-1\) 块。行列式为一迫使 \(-1\) 块的数目为偶数，可将它们两两写成 \(R(\pi)\)；\(1\) 块也两两写成 \(R(0)\)，奇数维只留一个 \(1\)。若对角化基的定向不合要求，反转其中一根基向量即可，至多把一个旋转角变号。于是可在特殊正交群内共轭入所列环面。沿各角的线性路径回到零，说明特殊正交群连通；它由正交方程定义为有界闭集，所以紧。

	为计算伴随作用，把标准实表示复化。第 \(j\) 个旋转平面有权 \(\varepsilon_j,-\varepsilon_j\)，奇数维还有一个零权方向。记复双线性标准配对为 \(B\)。映射
	\[
		u\wedge v\longmapsto\bigl(w\longmapsto B(u,w)v-B(v,w)u\bigr)
	\]
	给出 \(\Lambda^2\mathbb C^m\cong\mathfrak{so}(m,\mathbb C)\)，其中 \(m=2r\) 或 \(2r+1\)。在标准基上，它把 \(e_j\wedge e_k\) 送到基本反对称矩阵，因此是线性同构；保持 \(B\) 的矩阵与此映射交换，故还是表示的同构。外积中两权相加，于是不同旋转平面贡献 \(\pm\varepsilon_j\pm\varepsilon_k\)，奇数维的零权方向再贡献 \(\pm\varepsilon_j\)。同一平面的正负权相加为零，共有 \(r\) 个零权方向，恰好构成 \(\mathfrak t_{\mathbb C}\)。因此与 \(\mathfrak t\) 交换的 Lie 代数就是 \(\mathfrak t\)，任何更大的环面都会给出更大的交换子代数，矛盾。

	正规化子置换这些实旋转平面，并可在各平面上反转定向，分别对应置换 \(\theta_j\) 及改变其符号。置换两个二维平面的行列式为一，反转一个平面的行列式为负。在奇数维，可同时反转剩余的一维轴以恢复总行列式一，所以任意符号变化都可实现；偶数维没有这根轴，必须反转偶数个平面。对平面排列及其定向作用平凡的部分就是各平面内的旋转，模掉环面后不再有其他元素。这证明了两个 Weyl 群公式。
\end{proof}

奇数维与偶数维的上述根模型分别记作 \(B_r,D_r\)。\(\operatorname{SO}(2)\cong S^1\) 则是无根的交换群，不能把它当作有非零根的普通 \(D_r\) 例子。\(D_2\) 的四个根分成两组互相正交的正负根，这将在 \(\operatorname{SO}(4)\) 的双覆盖中得到解释。

\ProseParagraphBreak

本节的\term[jinxinqun]{紧辛群}[compact symplectic group]用 \(\operatorname{Sp}(r)\) 表示，它是保持四元数 Hermitian 内积的群，与保持实反对称双线性型的非紧群 \(\operatorname{Sp}(2r,\mathbb R)\) 不同。四元数代数 \(\mathbb H\) 有实基 \(1,\mathbf i,\mathbf j,\mathbf k\)，满足 \(\mathbf i^2=\mathbf j^2=\mathbf k^2=-1\)、\(\mathbf i\mathbf j=\mathbf k=-\mathbf j\mathbf i\)。共轭把虚部取负，且 \(\overline{ab}=\bar b\bar a\)、\(a\bar a=|a|^2\)，故非零元的逆是 \(\bar a/|a|^2\)。把向量写成右 \(\mathbb H\) 向量空间中的列，定义
\[
	\operatorname{Sp}(r)=\{A\in\operatorname{Mat}_{r\times r}(\mathbb H):A^*A=I\},\qquad A^*=\overline A^{\mathsf T}.
\]

\begin{proposition}[紧辛群的根模型]\label{g1:prop:symplectic-torus-roots}
	设 \(r\ge1\)，\(\mathbb H\) 为四元数体，固定四元数单位 \(\mathbf i\)，并令 \(\operatorname{Sp}(r)=\{A\in\operatorname{Mat}_{r\times r}(\mathbb H):A^*A=I\}\)，其中 \(A^*\) 表示四元数共轭转置。则 \(\operatorname{Sp}(r)\) 是实维数 \(r(2r+1)\) 的紧连通 Lie 群。其极大环面可取 \(\operatorname{diag}(e^{\mathbf i\theta_1},\ldots,e^{\mathbf i\theta_r})\)，整数格为 \(2\pi\mathbb Z^r\)。令 \(H(\theta)=\operatorname{diag}(\mathbf i\theta_1,\ldots,\mathbf i\theta_r)\)，用 \(\varepsilon_j(H(\theta))=\theta_j\) 表示坐标权，则
	\[
		R=\{\pm\varepsilon_j\pm\varepsilon_k:j<k\}\cup\{\pm2\varepsilon_j\},\qquad
		W\cong(\mathbb Z/2)^r\rtimes S_r.
	\]
\end{proposition}
\begin{proof}
	映射 \(A\mapsto A^*A\) 的值在四元数 Hermitian 矩阵空间中。在 \(A^*A=I\) 处，微分是 \(V\mapsto A^*V+V^*A\)；对任意 Hermitian 矩阵 \(S\)，取 \(V=AS/2\) 就得到 \(S\)。因而正则值定理给出嵌入子流形结构。Hermitian 矩阵有 \(r\) 个实对角元及 \(r(r-1)/2\) 个任意四元数上三角元，实维数为 \(2r^2-r\)，故群的维数为 \(4r^2-(2r^2-r)=r(2r+1)\)。乘法和求逆光滑，紧性来自闭性及每列长度为一。

	四元数 Gram--Schmidt 与实、复情形采用同一投影步骤：从 \(v\) 中减去 \(\sum_j u_j\langle u_j,v\rangle\)，再除以正实范数。右乘系数的次序不能交换。这说明群传递地作用于 \(S^{4r-1}\)，第一列为指定单位向量的稳定子同构于 \(\operatorname{Sp}(r-1)\)。\cref{g1:thm:homogeneous-space,g1:prop:transitive-action-model}把此球面识别为稳定子的商，并给出局部截面，故取第一列是局部平凡丛。\(\operatorname{Sp}(1)=S^3\) 连通；归纳假定纤维连通，沿球面路径用有限个局部平凡化提升，再在终点纤维内连接，就能把任一矩阵连到单位矩阵。因此各 \(\operatorname{Sp}(r)\) 都连通。

	写四元数矩阵为 \(Z+W\mathbf j\)，其中 \(Z,W\) 是复矩阵。代入 \(\mathbf jz=\bar z\mathbf j\) 可验证
	\[
		Z+W\mathbf j\longmapsto
		\begin{pmatrix}Z&W\\-\bar W&\bar Z\end{pmatrix}
	\]
	保持乘法及共轭转置。故 Lie 代数可写成
	\[
		\begin{pmatrix}A&B\\-\bar B&\bar A\end{pmatrix},\qquad A^*=-A,\quad B^{\mathsf T}=B.
	\]
	环面在这个复矩阵模型中为 \(\operatorname{diag}(e^{\mathrm{i}\theta_1},\ldots,e^{\mathrm{i}\theta_r},e^{-\mathrm{i}\theta_1},\ldots,e^{-\mathrm{i}\theta_r})\)。\(A\) 的非对角项给出 \(\varepsilon_j-\varepsilon_k\)，\(B\) 的非对角对称项给出 \(\varepsilon_j+\varepsilon_k\)，其对角项给出 \(2\varepsilon_j\)；共轭块给出各负根。零权空间只有 \(A\) 的纯虚对角元，故所列环面极大。

	置换四元数坐标实现任意 \(S_r\) 作用，在某一坐标左乘单位四元数 \(\mathbf j\)，则由 \(\mathbf j e^{\mathbf i\theta}\mathbf j^{-1}=e^{-\mathbf i\theta}\) 实现单个符号反转。这些矩阵属于紧辛群。反过来，正规化子在标准复表示中只能置换成对的权线 \(\varepsilon_j,-\varepsilon_j\)，所以作用必是带符号的置换；作用平凡的部分是环面。这就确定了 Weyl 群。
\end{proof}

这组根记作 \(C_r\)。\(B_r\) 与 \(C_r\) 有相同的反射超平面与 Weyl 群，但轴向根分别是 \(\varepsilon_j\) 和 \(2\varepsilon_j\)，它们与斜向根的相对长度不同。只画反射墙会丢掉这一差别；比较根模型还必须保留根向量与所用坐标尺度。

\subsection{低秩同构、中心及基本群}
\label{g1:subsec:1905:03}
根与 Weyl 群已经能控制许多共轭计算，但它们尚未决定中心与覆盖关系。我们先构造低维覆盖，再用球面上的标架丛计算基本群，避免把覆盖次数从根图中猜出来。

\begin{proposition}[四元数给出的低维群映射]\label{g1:prop:quaternion-low-rank}
	有光滑群同构 \(\operatorname{Sp}(1)\cong\operatorname{SU}(2)\)，其共同底流形为 \(S^3\)。另有满射 Lie 群同态
	\[
		\operatorname{Sp}(1)\longrightarrow\operatorname{SO}(3),\qquad
		\operatorname{Sp}(1)\times\operatorname{Sp}(1)\longrightarrow\operatorname{SO}(4),\qquad
		\operatorname{Sp}(2)\longrightarrow\operatorname{SO}(5),
	\]
	其核依次为 \(\{\pm1\}\)、\(\{(1,1),(-1,-1)\}\)、\(\{\pm I\}\)，所以均为二重覆盖。
\end{proposition}
\begin{proof}
	单位四元数 \(z+w\mathbf j\) 对应
	\[
		\begin{pmatrix}z&w\\-\bar w&\bar z\end{pmatrix},\qquad |z|^2+|w|^2=1.
	\]
	其列正交且行列式为一；反过来，特殊酉二阶矩阵的第二行由第一行及行列式一唯一确定为这一形式。上面已经核对乘法相容，所以得到光滑群同构及 \(S^3\) 模型。

	单位四元数在三维虚四元数空间上作用为 \(v\mapsto qvq^{-1}\)。共轭保持虚部与范数，且从单位元连续变化，故像在 \(\operatorname{SO}(3)\)。若作用恒等，\(q\) 与 \(\mathbf i,\mathbf j,\mathbf k\) 都交换，逐项比较便知 \(q\) 是实数，故核为 \(\{\pm1\}\)。源、靶同为三维，离散核使单位元微分单射，因而为同构。像是开子群，又因源紧而闭；靶连通，所以像为整个群。\cref{g1:prop:lie-homomorphism-local-isomorphism}于是给出二重覆盖。也可直接令 \(q=\cos t+u\sin t\)，其中 \(u\) 是单位虚四元数：共轭固定 \(u\)，在 \(u^\perp\) 上旋转 \(2t\)。这解释了\cref{g1:exm:su2-so3-lattices}中的两倍角。

	在四维实空间 \(\mathbb H\) 上令 \((a,b)\) 作用为 \(x\mapsto axb^{-1}\)。它保持范数，且源连通，像在 \(\operatorname{SO}(4)\)。若作用恒等，代入 \(x=1\) 得 \(a=b\)，再代入所有四元数得 \(a=b=\pm1\)。两群维数均为六，刚才的开闭像论证再次适用，得到第二个覆盖。

	最后令 \(V\) 为迹零的二阶四元数 Hermitian 矩阵空间：
	\[
		V=\left\{\begin{pmatrix}r&q\\\bar q&-r\end{pmatrix}:r\in\mathbb R,\ q\in\mathbb H\right\},\qquad
		\|X\|^2=\tfrac12\operatorname{Re}\operatorname{tr}(X^2)=r^2+|q|^2.
	\]
	它是五维实 Euclid 空间。\(A\in\operatorname{Sp}(2)\) 的作用为 \(X\mapsto AXA^*\)。四元数满足 \(\operatorname{Re}(ab)=\operatorname{Re}(ba)\)，按分量求和便得到实部迹的循环性；因而这个作用保持迹零条件及所写范数。源连通，像落在 \(\operatorname{SO}(V)\)。核中的 \(A\) 必须与 \(\operatorname{diag}(1,-1)\) 交换，所以 \(A=\operatorname{diag}(a,d)\)；再与所有非对角 \(X\) 交换，得 \(aq=qd\) 对所有 \(q\in\mathbb H\) 成立。取 \(q=1,\mathbf i,\mathbf j\) 得 \(a=d=\pm1\)。两群维数都是十，同样得到满射二重覆盖。
\end{proof}

\begin{lemma}[球面上的标架丛与基本群]\label{g1:lem:sphere-bundle-pi1}
	设 \(K\to E\xrightarrow{p}S^m\) 是光滑主 \(K\) 丛，\(K\) 为连通 Lie 群，\(m\ge3\)。则纤维包含 \(K\hookrightarrow E\) 诱导基本群同构。
\end{lemma}
\begin{proof}
	取分别去掉南、北极的开集 \(U,V\)，二者都微分同胚于 \(\mathbb R^m\)，交集同伦等价于 \(S^{m-1}\)。丛在 \(U,V\) 上各有全局截面。为说明这一点，可按\cref{g1:thm:principal-connection-existence}取主联络，在 Euclid 坐标中从原点向每点作径向路径，再从固定纤维点水平提升；\cref{g1:thm:principal-horizontal-lift}保证提升在整个区间存在并光滑依赖终点，因此给出所需截面。这只用主联络沿紧区间的提升，不要求底空间完备。

	于是 \(p^{-1}(U)\) 与 \(p^{-1}(V)\) 都是可缩开集与 \(K\) 的乘积；交集用任一平凡化写成 \((U\cap V)\times K\)。因 \(m-1\ge2\)，\cref{g1:prop:sphere-simply-connected}说明 \(U\cap V\) 单连通。乘积中的路径与同伦逐分量进行，所以这三个空间的基本群都通过同一基点纤维与 \(\pi_1(K)\) 同构。两条包含映射在该纤维上就是恒等，故在这些识别下也都是恒等。\cref{g1:thm:van-kampen}将两个 \(\pi_1(K)\) 沿共同的同一 \(\pi_1(K)\) 粘合，所得群仍为 \(\pi_1(K)\)，且同构正由纤维包含诱导。
\end{proof}

\begin{proposition}[经典紧群的中心与基本群]\label{g1:prop:classical-centers-pi1}
	下表列出相应连通经典紧群的中心与基本群。其中 \(\mu_n=\{z\in S^1:z^n=1\}\)，\(r\ge1\)，并单独保留低维旋转群。
	\[
		\begin{array}{c|c|c}
		G&Z(G)&\pi_1(G)\\ \hline
		\operatorname{U}(n),\ n\ge1&S^1I&\mathbb Z\\
		\operatorname{SU}(n),\ n\ge2&\mu_nI&0\\
		\operatorname{Sp}(r)&\{\pm I\}&0\\
		\operatorname{SO}(2)&\operatorname{SO}(2)&\mathbb Z\\
		\operatorname{SO}(2r+1),\ r\ge1&\{I\}&\mathbb Z/2\\
		\operatorname{SO}(2r),\ r\ge2&\{\pm I\}&\mathbb Z/2.
		\end{array}
	\]
\end{proposition}
\begin{proof}
	酉或特殊酉群的中心元必须与其对角环面交换，所以为对角矩阵；再与混合任意两个坐标的特殊酉二阶块交换，迫使所有对角元相同。行列式条件分别给出 \(S^1I\) 与 \(\mu_nI\)。紧辛群中，先与复对角环面交换得到复对角矩阵，再与各坐标上的 \(\mathbf j\) 交换，迫使对角元为实数；置换坐标使它们相同，故只有 \(\pm I\)。对于 \(m\ge3\)，若实线性变换与所有坐标平面旋转交换，则对每个单位向量 \(v\)，它必须保持固定 \(v\) 的旋转子群的共同不动空间 \(\mathbb Rv\)。因此每条直线不变，由线性性是标量变换。正交性及行列式一分别给出奇数维的 \(I\) 和偶数维的 \(\pm I\)。\(\operatorname{SO}(2)\) 本来交换，须另行处理。

	取第一列给出主丛
	\[
		\operatorname{SU}(n-1)\longrightarrow\operatorname{SU}(n)\longrightarrow S^{2n-1},\qquad
		\operatorname{Sp}(r-1)\longrightarrow\operatorname{Sp}(r)\longrightarrow S^{4r-1}.
	\]
	第一条中的传递性由酉 Gram--Schmidt 得到：选定第一列后，用其余列之一调整行列式而不改变第一列。稳定子恰为低一阶的特殊酉群。第二条的传递性和稳定子已经在\cref{g1:prop:symplectic-torus-roots}中证明。它们均由齐性空间的局部截面构成主丛，符合\cref{g1:exm:homogeneous-principal-bundle}的条件。以 \(\operatorname{SU}(1)=\operatorname{Sp}(0)=\{1\}\) 为起点，反复应用上一引理，得到两类群的基本群都为零。\(\operatorname{SU}(2)\cong S^3\) 还给出独立的起始核验。

	映射
	\[
		S^1\times\operatorname{SU}(n)\longrightarrow\operatorname{U}(n),\qquad
		(z,A)\longmapsto\operatorname{diag}(z,1,\ldots,1)A
	\]
	是微分同胚：逆映射先取 \(z=\det U\)，再取 \(A=\operatorname{diag}(z^{-1},1,\ldots,1)U\)。它通常不是直积群同构，但已足够说明 \(\pi_1(\operatorname{U}(n))\cong\pi_1(S^1)=\mathbb Z\)，且行列式诱导此同构。

	\(S^3\to\operatorname{SO}(3)\) 是前一命题给出的二重覆盖；\(S^3\) 单连通，所以\cref{g1:thm:normal-covering-deck}将基本群识别为二元覆盖变换群，给出 \(\pi_1(\operatorname{SO}(3))\cong\mathbb Z/2\)。对 \(n\ge4\)，第一列又给出主丛
	\[
		\operatorname{SO}(n-1)\longrightarrow\operatorname{SO}(n)\longrightarrow S^{n-1}.
	\]
	正交标架可在不改变第一列的条件下调整定向，故作用传递，稳定子为 \(\operatorname{SO}(n-1)\)。球面维数至少三，上一引理说明包含映射诱导基本群同构，归纳便得全部 \(n\ge3\) 的结论。\(\operatorname{SO}(2)\) 的一整周旋转生成 \(\mathbb Z\)，进入 \(\operatorname{SO}(3)\) 后只保留其奇偶性。
\end{proof}

至此，\(\operatorname{Sp}(1)\times\operatorname{Sp}(1)\to\operatorname{SO}(4)\) 与 \(\operatorname{Sp}(2)\to\operatorname{SO}(5)\) 的源群也已证明单连通，因而是相应的万有覆盖。微分给出 \(\mathfrak{so}(4)\cong\mathfrak{su}(2)\oplus\mathfrak{su}(2)\) 及 \(\mathfrak{so}(5)\cong\mathfrak{sp}(2)\)，但群本身仍须除以所列二元中心核。根模型中的 \(D_2\) 分裂及 \(B_2,C_2\) 的低秩对应因此各有具体的矩阵实现。

\ProseParagraphBreak

这些算例显示三种资料各自回答不同问题。根给出共轭轨道的无穷小方向，Weyl 群给出同一环面内的有限识别，整数格与中心商则决定哪些环路真正闭合以及覆盖会怎样改变群。后一类信息不能从 Lie 代数同构中自动恢复。


\begin{chaptersummary}
	紧群上的归一化不变积分把任意内积平均为伴随不变内积，左平移后得到双不变度量。径向长度比较与紧性共同保证任意两点之间存在最短测地线；在双不变度量中，从单位元出发的测地线是一参数子群。因此紧连通 Lie 群的每个元素都可以写成 \(\exp X\)。

	\ProseParagraphBreak

	一参数子群的闭包是紧连通交换 Lie 群，而这样的群具有 \(\mathfrak t/\Lambda\) 的形式，其中 \(\Lambda\) 是满秩格。维数的有界性使每个环面包含在极大环面中。于是每个群元素属于某个极大环面；伴随轨道上的极值论证进一步证明所有极大环面共轭。固定一个极大环面后，群内的共轭问题便集中到这个环面。

	\ProseParagraphBreak

	中心化子说明哪些方向仍与给定对象交换，正规化子还允许重新安排环面内部的方向。有限商 \(W=N_G(T)/Z_G(T)\) 描述后一种作用。伴随表示的权分解把非交换方向组织为根空间，根超平面分隔正则方向，根反射由群内变换实现。群元素的正则条件涉及 \(e^{\mathrm{i}\alpha(H)}\)，Lie 代数元素的正则条件则涉及 \(\alpha(H)\)，两者不能直接混同。

	\ProseParagraphBreak

	酉群给出差根 \(\varepsilon_j-\varepsilon_k\)，旋转群给出 \(B_r,D_r\) 模型，紧辛群给出 \(C_r\) 模型。矩阵计算同时确定整数格、Weyl 群和中心。\(\operatorname{SU}(2)\to\operatorname{SO}(3)\) 等低维覆盖说明：微分同构保留 Lie 代数，却可能丢失离散中心核。特征标能否下降、环路绕行多少次及群是否单连通，都需要进一步保留这些整体信息。
\end{chaptersummary}

\BookExercises

以下习题从平均与最短路径开始，逐步转到环面的闭包、共轭、根与经典群计算。涉及指数时分别检查 Lie 代数中的等式和模整数格后的群等式；涉及共轭时明确允许的共轭群。

\begin{exercise}\label{g1:ex:unitary-conjugation-average}
	对任意复矩阵 \(A\in M_n(\mathbb C)\)，令
	\[
		P(A)=\int_{\operatorname{U}(n)}UAU^{-1}\,\mathrm{d}\mu(U).
	\]
	证明 \(P(A)=\frac{\operatorname{tr}A}{n}I\)。进而求 \(\int_{\operatorname{U}(n)}|U_{11}|^2\,\mathrm{d}\mu(U)\)，并在 \(n>1\) 时举例说明，共轭平均不保持矩阵的秩。
\end{exercise}
\begin{solution}
	对任意 \(V\in\operatorname{U}(n)\)，左不变换元 \(U\mapsto VU\) 给出 \(VP(A)V^{-1}=P(A)\)。若一个矩阵与所有对角酉矩阵交换，则其每个非对角元均为零：独立改变第 \(j\) 与第 \(k\) 个对角相位，会把该位置乘以任意单位复数，而不变性要求原值不变。故 \(P(A)\) 对角。再与所有置换矩阵交换，得到各对角元相同，所以 \(P(A)=cI\)。迹与有限分量积分交换，且 \(\operatorname{tr}(UAU^{-1})=\operatorname{tr}A\)，于是 \(nc=\operatorname{tr}A\)，证明公式。

	令 \(A=E_{11}\)。因 \(U^{-1}=U^*\)，矩阵 \(UE_{11}U^*\) 的第 \((1,1)\) 项为 \(|U_{11}|^2\)，而平均为 \(I/n\)，故所求积分是 \(1/n\)。\(E_{11}\) 的秩为 \(1\)，但 \(P(E_{11})=I/n\) 的秩为 \(n\)；当 \(n>1\) 时二者不同。
\end{solution}

\begin{exercise}\label{g1:ex:bi-invariant-curvature}
	设有限维实 Lie 群 \(G\) 配备双不变度量。采用本书的曲率约定，证明对单位元切向量 \(X,Y,Z\)，有
	\[
		R(X,Y)Z=-\tfrac14[[X,Y],Z].
	\]
	当 \(X,Y\) 线性无关时，由它们张成平面的\term[jiemianqulv]{截面曲率}[sectional curvature]定义为
	\[
		K(X,Y)=\frac{\langle R(X,Y)Y,X\rangle}{|X|^2|Y|^2-\langle X,Y\rangle^2}.
	\]
	求此处的截面曲率，证明其非负，并证明该度量平坦当且仅当 \(\mathfrak g\) 交换。最后说明没有连通性时不能由平坦推出 \(G\) 交换。
\end{exercise}
\begin{solution}
	用左不变场延拓 \(X,Y,Z\)。由\eqref{g1:eq:bi-invariant-connection}和曲率定义，
	\[
		R(X,Y)Z
		=\tfrac14[X,[Y,Z]]-\tfrac14[Y,[X,Z]]-\tfrac12[[X,Y],Z].
	\]
	Jacobi 恒等式使前两项之和为 \(\tfrac14[[X,Y],Z]\)，故得到所求公式。记 \(A=[X,Y]\)。伴随不变内积的斜伴随恒等式给出
	\[
		\langle[A,Y],X\rangle
		=-\langle[Y,A],X\rangle
		=\langle A,[Y,X]\rangle=-|A|^2.
	\]
	所以
	\[
		K(X,Y)=\frac{|[X,Y]|^2}
		{4\bigl(|X|^2|Y|^2-\langle X,Y\rangle^2\bigr)}\ge0.
	\]
	若 \(\mathfrak g\) 交换，曲率公式直接给出 \(R=0\)。反之，若 \(R=0\)，对任意线性无关的 \(X,Y\) 有 \([X,Y]=0\)；线性相关时括号也为零，故整个 Lie 代数交换。取非交换有限群 \(F\)，并令 \(G=S^1\times F\)。它的每个分支都是标准圆周，可由 \(\mathfrak g=\mathbb R\) 上的标准内积得到双不变平坦度量，但群乘法包含非交换的 \(F\)，所以 \(G\) 不交换。
\end{solution}

\begin{exercise}\label{g1:ex:compact-group-power-map}
	设 \(G\) 为有限维实紧连通 Lie 群，\(m\ge1\) 为整数。证明幂映射 \(P_m:G\to G\)、\(P_m(g)=g^m\)，满射。用 \(\operatorname{SU}(2)\) 中的具体矩阵说明 \(P_2\) 未必是群同态，并用 \(\operatorname{O}(2)\) 说明去掉连通性后 \(P_2\) 未必满射。
\end{exercise}
\begin{solution}
	给定 \(g\in G\)，\cref{g1:thm:compact-exp-surjective}给出 \(X\in\mathfrak g\) 使 \(g=\exp_GX\)。令 \(h=\exp_G(X/m)\)，单参数子群规律给出 \(h^m=\exp_GX=g\)，因此满射。

	在 \(\operatorname{SU}(2)\) 中取
	\[
		A=\begin{pmatrix}\mathrm{i}&0\\0&-\mathrm{i}\end{pmatrix},
		\qquad B=\begin{pmatrix}0&1\\-1&0\end{pmatrix}.
	\]
	二者均酉且行列式为 \(1\)，并满足 \(A^2=B^2=-I\)、\(AB=-BA\)。所以 \((AB)^2=-A^2B^2=-I\)，但 \(A^2B^2=I\)，幂映射不保持乘法。在 \(\operatorname{O}(2)\) 中，任意 \(Q\) 均有 \(\det(Q^2)=1\)，故行列式为 \(-1\) 的反射没有平方根。连通性在满射命题中不能省略。
\end{solution}

\begin{exercise}\label{g1:ex:flat-torus-shortest-logarithms}
	在 \(T^r=\mathbb R^r/\mathbb Z^r\) 上取由正定对称矩阵 \(Q\) 给出的常系数度量，记 \(|v|_Q^2=v^{\mathsf T}Qv\)。证明
	\[
		d(0,x+\mathbb Z^r)=\min_{z\in\mathbb Z^r}|x-z|_Q.
	\]
	当 \(r=2\)、\(Q=\left(\begin{smallmatrix}1&9/10\\9/10&1\end{smallmatrix}\right)\)、\(x=(49/100,49/100)\) 时，求所有最短测地线的初速度，并说明为什么逐坐标选取 \([-1/2,1/2]\) 中的代表元未必得到最短对数。
\end{exercise}
\begin{solution}
	商映射是局部等距映射，且指数映射为 \(v\mapsto v+\mathbb Z^r\)。\cref{g1:thm:compact-exp-surjective}给出一条初速度 \(v\) 的最短测地线；其终点条件为 \(v=x-z\)，其中 \(z\in\mathbb Z^r\)，故 \(d(0,x+\mathbb Z^r)=|x-z|_Q\)。反过来，每个这样的 \(x-z\) 都给出一条连接两点的直线投影，其长度为 \(|x-z|_Q\)，所以距离不大于任何一个候选值，证明最小值公式。也可直接从 \(Q\) 正定看出候选值在 \(|z|\to\infty\) 时趋于无穷，因此最小值确能取得。

	写 \(z=(m,n)\)，并令 \(u=49/100-m\)、\(v=49/100-n\)。则
	\[
		|(u,v)|_Q^2
		=\frac{19}{20}(u+v)^2+\frac1{20}(u-v)^2
		=\frac{19}{20}\bigl(49/50-m-n\bigr)^2
		+\frac1{20}(n-m)^2.
	\]
	若 \(m+n\) 为偶数，则 \(|49/50-m-n|\ge49/50\)，故平方长度至少为 \(45619/50000\)，在 \(m=n=0\) 时取得。若 \(m+n\) 为奇数，则 \(|49/50-m-n|\ge1/50\) 且 \(|n-m|\ge1\)，故平方长度至少为 \(2519/50000\)。等号要求 \(m+n=1\)、\(|n-m|=1\)，恰好给出 \((m,n)=(1,0)\) 或 \((0,1)\)。因此最短测地线的初速度恰有两个：
	\[
		(-51/100,49/100),\qquad(49/100,-51/100).
	\]
	原代表元 \((49/100,49/100)\) 的每个坐标都在指定半开或闭区间内，却具有较大长度。原因是 \(Q\) 的交叉项使两个坐标方向并不正交，逐坐标最小化不能代替按完整二次型最小化。
\end{solution}

\begin{exercise}\label{g1:ex:torus-integer-homomorphism}
	把 \(n\) 维环面写为 \(T^n=\mathbb R^n/\mathbb Z^n\)。证明每个光滑群同态 \(F:T^n\to T^m\) 都唯一写成
	\[
		F(x+\mathbb Z^n)=Ax+\mathbb Z^m,
	\]
	其中 \(A\) 是一个 \(m\times n\) 整数矩阵。进而证明 \(T^n\) 的 Lie 群自同构恰对应于 \(\operatorname{GL}(n,\mathbb Z)\)，并说明一个可逆实矩阵未必诱导环面自同构。
\end{exercise}
\begin{solution}
	商环面的 Lie 代数为 \(\mathbb R^n\)，其指数映射为 \(x\mapsto x+\mathbb Z^n\)。令 \(A=(\mathrm{d}F)_0\)，指数自然性给出
	\[
		F(x+\mathbb Z^n)=Ax+\mathbb Z^m.
	\]
	若 \(z\in\mathbb Z^n\)，左侧为单位元，故 \(Az\in\mathbb Z^m\)。对每个标准基向量应用此条件，可知 \(A\) 的每一列为整数列。反之，整数矩阵满足 \(A\mathbb Z^n\subseteq\mathbb Z^m\)，所以公式与代表元无关，且由线性映射下降成光滑群同态。其微分就是 \(A\)，因此矩阵唯一。

	若 \(F\) 为 Lie 群自同构，\(F^{-1}\) 同样由整数矩阵 \(B\) 给出，链式法则使 \(AB=BA=I\)。因而 \(\det A\)、\(\det B\) 都是整数且乘积为 \(1\)，故 \(\det A=\pm1\)。反之，整数矩阵若行列式为 \(\pm1\)，伴随矩阵公式给出整数逆矩阵，遂诱导逆同态。故自同构群恰为 \(\operatorname{GL}(n,\mathbb Z)\)。例如一维的 \(A=\sqrt2\) 不保持整数格，连良定义的商映射也不给出；\(A=2\) 虽然可逆且保持整数格，但诱导的圆周映射有两个原像，因而不是自同构。
\end{solution}

\begin{exercise}\label{g1:ex:three-torus-flow-closure}
	在 \(T^3=(S^1)^3\) 中考虑单参数子群
	\[
		c(t)=(e^{\mathrm{i}t},e^{\mathrm{i}\sqrt2t},e^{\mathrm{i}(1+\sqrt2)t}).
	\]
	求它的像的闭包及闭包的 Lie 代数，并比较该 Lie 代数与 \(c'(0)\) 张成的直线。
\end{exercise}
\begin{solution}
	每个 \(c(t)\) 都满足 \(z_3=z_1z_2\)，所以闭包包含于闭子群
	\[
		A=\{(z_1,z_2,z_1z_2):z_1,z_2\in S^1\}.
	\]
	映射 \(j:T^2\to A\)、\(j(z_1,z_2)=(z_1,z_2,z_1z_2)\) 是 Lie 群同构，逆为前两个坐标的投影。由\cref{g1:exm:irrational-torus-subgroup}，\(t\mapsto(e^{\mathrm{i}t},e^{\mathrm{i}\sqrt2t})\) 在 \(T^2\) 中稠密，因为 \(\sqrt2\) 无理。将其送入同胚 \(j\)，可知 \(c(\mathbb R)\) 在 \(A\) 中稠密。因此所求闭包恰为 \(A\)。

	在单位元处微分 \(j\)，得到
	\[
		\operatorname{Lie}(A)=\{(\mathrm{i}u,\mathrm{i}v,\mathrm{i}(u+v)):u,v\in\mathbb R\}.
	\]
	这是二维实平面，而 \(c'(0)=\mathrm{i}(1,\sqrt2,1+\sqrt2)\) 只张成其中一条直线。闭包增加了一维的无穷小方向，尽管所有原像点都由同一个实参数产生。
\end{solution}

\begin{exercise}\label{g1:ex:commuting-elements-no-common-torus}
	在 \(\operatorname{SO}(3)\) 中令
	\[
		A=\operatorname{diag}(1,-1,-1),\qquad
		B=\operatorname{diag}(-1,1,-1).
	\]
	证明 \(A\) 与 \(B\) 交换，但不存在同时包含二者的环面。说明这与“每个元素属于某个极大环面”并不矛盾。
\end{exercise}
\begin{solution}
	两矩阵都正交且行列式为 \(1\)，作为对角矩阵彼此交换。它们满足 \(A^2=B^2=I\)，且 \(I,A,B,AB\) 是四个不同元素。

	先看绕第三轴旋转的标准圆周 \(T\subseteq\operatorname{SO}(3)\)。若 \(C\) 与 \(T\) 的所有元素交换，便保持 \(T\) 的公共不动轴 \(\mathbb Re_3\) 及其正交平面。它在平面上的限制与所有平面旋转交换。写一个实 \(2\times2\) 矩阵并令它与 \(\begin{psmallmatrix}0&-1\\1&0\end{psmallmatrix}\) 交换，可得它形如 \(\begin{psmallmatrix}a&-b\\b&a\end{psmallmatrix}\)；正交性给出 \(a^2+b^2=1\)，故这是行列式为 \(1\) 的平面旋转。\(C\) 在第三轴上的因子由 \(\det C=1\) 也等于 \(1\)。所以 \(Z_{\operatorname{SO}(3)}(T)=T\)，任何包含 \(T\) 的环面都只能等于 \(T\)。这证明 \(T\) 极大。

	\(\operatorname{SO}(3)\) 连通：每个三维正交行列式为 \(1\) 的矩阵都有特征值 \(1\)，并在其正交平面上是旋转，连续缩小旋转角便给出到单位元的道路。因此极大环面共轭定理适用，所有极大环面都是圆周，任意环面又包含于一个极大环面。如果一个环面同时包含 \(A,B\)，相应圆周便至少有四个满足 \(x^2=e\) 的元素。然而 \(S^1\) 中方程 \(z^2=1\) 只有 \(1,-1\) 两解，矛盾。逐个元素的存在定理允许为不同元素选择不同环面，并没有给出有限个交换元素的共同环面。
\end{solution}

\begin{exercise}\label{g1:ex:center-intersection-maximal-tori}
	设 \(G\) 为紧连通 Lie 群。证明其中心恰为全部极大环面的交：
	\[
		Z(G)=\bigcap_{T\text{ 为极大环面}}T.
	\]
	由此说明，即使 \(G\) 连通，这个交也未必连通，并给出具体例子。
\end{exercise}
\begin{solution}
	若 \(z\in Z(G)\)，则 \(z\) 与任一极大环面 \(T\) 的所有元素交换。由\cref{g1:prop:maximal-torus-centralizer}，\(z\in Z_G(T)=T\)，所以中心包含于所列交集。

	反过来，若 \(z\) 属于每个极大环面，任取 \(g\in G\)。由\cref{g1:thm:element-in-maximal-torus}，可选包含 \(g\) 的极大环面 \(T_g\)。假设也给出 \(z\in T_g\)，由交换性得 \(zg=gz\)。\(g\) 任意，故 \(z\in Z(G)\)。

	例如 \(\operatorname{SU}(2)\) 的中心为 \(\{I,-I\}\)。可直接核对：与全部对角矩阵 \(\operatorname{diag}(e^{\mathrm{i}t},e^{-\mathrm{i}t})\) 交换的矩阵必须对角；这样的矩阵若再与 \(\begin{psmallmatrix}0&1\\-1&0\end{psmallmatrix}\) 交换，两对角元必须相等，行列式为 \(1\) 便给出 \(\pm I\)。反向这两个标量矩阵确实与全群交换。\(\operatorname{SU}(2)\) 同胚于 \(S^3\)，是连通群，而两个不同点组成的中心不连通。
\end{solution}

\begin{exercise}\label{g1:ex:compact-power-map}
	设 \(G\) 为紧连通 Lie 群，\(m\) 为正整数。证明一个 \(r\) 维环面中每个元素恰有 \(m^r\) 个 \(m\) 次方根。一般紧连通群中的根数是否必然有限？用 \(\operatorname{SO}(3)\) 中的平方根说明。
\end{exercise}
\begin{solution}
	在 \(T\cong\mathbb R^r/\mathbb Z^r\) 中，求 \(x\) 的 \(m\) 次方根就是求 \(my=x\pmod{\mathbb Z^r}\)。固定一个代表元 \(x\in\mathbb R^r\)，全部解为
	\[
		y=\frac{x+k}{m}+\mathbb Z^r,
		\qquad k\in\{0,1,\ldots,m-1\}^r.
	\]
	任意整数向量可模 \(m\) 化成上述唯一代表；不同代表的差除以 \(m\) 不在 \(\mathbb Z^r\)，故这些解彼此不同，恰有 \(m^r\) 个。

	全群中的根数却可以无限。\(\operatorname{SO}(3)\) 中，绕任意一条过原点的轴旋转 \(\pi\) 都满足 \(R^2=I\)，不同的无向轴给出不同矩阵。事实上若 \(R^2=I\)，则 \(R^{-1}=R\) 且正交性给出 \(R^{\mathsf T}=R\)。实对称谱分解使特征值全为 \(\pm1\)，再由行列式为 \(1\)，可知 \(R=I\)，或 \(R\) 在一维空间上为 \(1\)、其正交平面上为 \(-1\)。后一种恰是绕该直线的半周旋转。因此 \(I\) 的平方根集合由孤立点 \(I\) 及一个以 \(\mathbb{RP}^2\) 为参数的无限族组成；不同环面中的根不能按单个环面的有限计数合并。
\end{solution}

\begin{exercise}\label{g1:ex:root-height-a2}
	在 \(\operatorname{SU}(3)\) 的对角子代数中写 \(H=\mathrm{i}\operatorname{diag}(\theta_1,\theta_2,\theta_3)\)，其中 \(\theta_1+\theta_2+\theta_3=0\)。设 \(\alpha=\varepsilon_1-\varepsilon_2\)、\(\beta=\varepsilon_2-\varepsilon_3\)。用 \(B(X,Y)=-\operatorname{tr}(XY)\) 计算两个余根和对应反射；描述含 \(H_0=\mathrm{i}\operatorname{diag}(2,0,-2)\) 的 Weyl 室，并把 \(\mathrm{i}\operatorname{diag}(-3,2,1)\) 用反射移到该室。
\end{exercise}
\begin{solution}
	在角度坐标中，内积是 \(\mathbb R^3\) 的普通内积在和为零平面上的限制。因此
	\[
		A_\alpha=(1,-1,0),\quad A_\beta=(0,1,-1),\quad
		(\alpha,\alpha)=(\beta,\beta)=2,
	\]
	所以 \(\alpha^\vee=(1,-1,0)\)、\(\beta^\vee=(0,1,-1)\)。反射公式给出
	\[
		s_\alpha(\theta_1,\theta_2,\theta_3)=(\theta_2,\theta_1,\theta_3),\qquad
		s_\beta(\theta_1,\theta_2,\theta_3)=(\theta_1,\theta_3,\theta_2).
	\]
	六根为 \(\pm\alpha,\pm\beta,\pm(\alpha+\beta)\)。在 \(H_0\) 上三个所列正根都取正值，故相应室由 \(\theta_1>\theta_2>\theta_3\)、\(\sum\theta_j=0\) 定义。角度向量 \((-3,2,1)\) 先经 \(s_\alpha\) 变为 \((2,-3,1)\)，再经 \(s_\beta\) 变为 \((2,1,-3)\)，后者属于所求室。整个过程保持角度和与 Euclid 范数，并对应群内共轭。
\end{solution}

\begin{exercise}\label{g1:ex:b2-coroots}
	在 Euclid 平面 \(\mathbb R^2\) 中考虑根集
	\[
		R=\{\pm\varepsilon_1,\pm\varepsilon_2,
		\pm(\varepsilon_1+\varepsilon_2),\pm(\varepsilon_1-\varepsilon_2)\}.
	\]
	计算 \(Q,Q^\vee,P,P^\vee\)，并求 \(P/Q\) 与 \(P^\vee/Q^\vee\)。说明短根与长根取余根以后长度次序如何改变。
\end{exercise}
\begin{solution}
	通过标准内积识别向量与协向量。短根 \(\pm\varepsilon_j\) 的余根为 \(\pm2e_j\)；长根 \(\pm(\varepsilon_1\pm\varepsilon_2)\) 的余根为 \(\pm(e_1\pm e_2)\)。因此
	\[
		Q=\mathbb Z^2,\qquad
		Q^\vee=\{(m,n)\in\mathbb Z^2:m+n\text{ 为偶数}\}.
	\]
	后一个等式可由生成元 \((1,1),(1,-1)\) 看出：它们的整数线性组合恰好使两坐标同奇偶，而 \((2,0),(0,2)\) 已在其中。余权格由根的整数取值定义，由于 \(\varepsilon_1,\varepsilon_2\) 本身都是根，得到 \(P^\vee=\mathbb Z^2\)。

	权格要求 \(2a,2b,a+b,a-b\) 都是整数，等价于
	\[
		P=\mathbb Z^2\ \cup\ \bigl((\tfrac12,\tfrac12)+\mathbb Z^2\bigr).
	\]
	所以 \(P/Q\) 由半整数坐标的类生成，阶为二；\(P^\vee/Q^\vee\) 由坐标和的奇偶分类，也同构于 \(\mathbb Z/2\mathbb Z\)。原短根长为一、原长根长为 \(\sqrt2\)，余根相应长为二、\(\sqrt2\)，故长短次序反转。它解释了同一反射超平面排列仍能配上不同的根长数据。
\end{solution}

\begin{exercise}\label{g1:ex:regular-so3-disconnected-centralizer}
	令 \(g=\operatorname{diag}(-1,-1,1)\in\operatorname{SO}(3)\)，即绕第三坐标轴旋转 \(\pi\)。求 \(Z_{\operatorname{SO}(3)}(g)\) 及其单位分支，证明 \(g\) 是正则群元素，但其中心化子不是极大环面。
\end{exercise}
\begin{solution}
	与 \(g\) 交换的矩阵必须保持它的两个实特征空间 \(E_-\) 与 \(E_+\)，所以具有分块形式 \(\operatorname{diag}(A,\epsilon)\)，其中 \(A\in\operatorname{O}(2)\)、\(\epsilon\in\{1,-1\}\)。行列式等于一恰好要求 \(\epsilon=\det A\)，故
	\[
		Z_{\operatorname{SO}(3)}(g)
		=\{\operatorname{diag}(A,\det A):A\in\operatorname{O}(2)\}
		\cong\operatorname{O}(2).
	\]
	其单位分支为 \(\operatorname{diag}(\operatorname{SO}(2),1)\)，就是绕第三坐标轴的旋转环面 \(T\)。这个环面极大：与全部平面旋转交换的三维特殊正交矩阵必须保留轴线；在平面上只有旋转与全部旋转交换，行列式约束随后使轴线方向取正。故 \(Z_{\operatorname{SO}(3)}(T)=T\)，任何包含 \(T\) 的环面都只能是 \(T\)。于是群的秩为一，\(g\) 的中心化子也为一维，满足正则性定义。但是 \(\operatorname{O}(2)\) 有两个连通分支，例如 \(\operatorname{diag}(1,-1,-1)\) 在非单位分支中。因此整个中心化子严格大于 \(T\)。
\end{solution}

\begin{exercise}\label{g1:ex:su3-central-quotient-lattice}
	取 \(G=\operatorname{SU}(3)\) 及其中心子群 \(C=\{\zeta I:\zeta^3=1\}\)。用角度和为零的坐标描述 \(G/C\) 的周期格，并判定 \(\operatorname{SU}(3)\) 环面特征标
	\(\chi_m(\theta)=\exp(\mathrm{i}\sum_{j=1}^3m_j\theta_j)\)、\(m_j\in\mathbb Z\)，何时下降到商环面。由此求商环面的特征标格在原特征标格中的指数。
\end{exercise}
\begin{solution}
	指数落在 \(C\) 中，当且仅当三个相位相同且为三次单位根。因此周期格是
	\[
		\Lambda'=
		\{2\pi(m_1+k/3,m_2+k/3,m_3+k/3):
		m_j,k\in\mathbb Z,\ \sum_jm_j=-k\}.
	\]
	其中 \(k\) 改变三的倍数不改变相应陪集，所以原来的格 \(\Lambda=2\pi\{m\in\mathbb Z^3:\sum m_j=0\}\) 在 \(\Lambda'\) 中的指数为三。直接地，指数映射给出 \(\Lambda'/\Lambda\cong C\)。

	在和为零的平面上，两个整数三元组 \(m,m'\) 给出同一线性泛函，当且仅当 \(m-m'\) 是 \((1,1,1)\) 的整数倍。每个环面特征标都可如此表示：格基 \(2\pi(1,0,-1)\)、\(2\pi(0,1,-1)\) 上的取值除以 \(2\pi\) 是任意两个整数，可取 \(m_3=0\) 实现。对 \(\zeta I\)，有 \(\chi_m(\zeta I)=\zeta^{m_1+m_2+m_3}\)。故下降条件是 \(m_1+m_2+m_3\equiv0\pmod3\)；换一个表示三元组会给这个和加上三的倍数，条件不变。坐标和模三给出原特征标格到 \(\mathbb Z/3\mathbb Z\) 的满同态，其核就是可下降的格，故指数为三。
\end{solution}

\begin{exercise}\label{g1:ex:root-data-central-torus}
	设 \(K\) 为紧连通 Lie 群，\(T_K\) 为其极大环面，\(r\ge1\)。求 \(G=K\times(S^1)^r\) 的一个极大环面、Weyl 群、根集、周期格及特征标格，并描述它的 Weyl 室。说明单凭根集为何不能记录新增的连续中心方向。
\end{exercise}
\begin{solution}
	令 \(T=T_K\times(S^1)^r\)。若一个环面包含 \(T\)，它在 \(K\) 中的投影是包含 \(T_K\) 的紧连通交换子群，故由极大性等于 \(T_K\)；于是原环面已包含于 \(T\)，证明 \(T\) 极大。第二因子为中心，故
	\[
		N_G(T)=N_K(T_K)\times(S^1)^r,\qquad
		W(G,T)\cong W(K,T_K).
	\]
	Lie 代数为 \(\mathfrak k\oplus\mathbb R^r\)，第二个直和项的括号恒零，因此根恰为 \((\alpha,0)\)、\(\alpha\in R_K\)，各根空间来自原来的 \(\mathfrak k_\alpha\)。取标准角度坐标后，
	\[
		\Lambda_T=\Lambda_{T_K}\oplus2\pi\mathbb Z^r,\qquad
		X^*(T)=X^*(T_K)\oplus\mathbb Z^r.
	\]
	全部根超平面为 \(\ker\alpha\times\mathbb R^r\)，所以每个室是 \(C_K\times\mathbb R^r\)。根格始终为 \(Q_K\oplus\{0\}\)，新增维数全部是零权方向。若只保留根张成空间及其中的根集，各个 \(r\) 得到同一份数据；要区分这些群，至少还须记录完整的环面空间及其周期格。
\end{solution}

\begin{exercise}\label{g1:ex:unitary-multiplicity-orbit}
	设 \(\lambda,\mu,\nu\in S^1\) 两两不同，且 \(g=\operatorname{diag}(\lambda,\lambda,\mu,\nu)\in\operatorname{U}(4)\)。求 \(Z_{\operatorname{U}(4)}(g)\) 及 \(g\) 的共轭轨道维数。若再假设 \(\lambda^2\mu\nu=1\)，求特殊酉群中的中心化子及共轭轨道维数。说明为什么两种轨道的维数相同，而稳定子并不相同。
\end{exercise}
\begin{solution}
	与 \(g\) 交换的线性映射必须保持三个特征子空间 \(\mathbb C^2\)、\(\mathbb C\)、\(\mathbb C\)，在各子空间上又可任取酉变换。因此
	\[
		Z_{\operatorname{U}(4)}(g)=\operatorname{U}(2)\times\operatorname{U}(1)\times\operatorname{U}(1),
	\]
	实维数为 \(4+1+1=6\)。由轨道的齐性空间模型，共轭轨道维数为 \(16-6=10\)。若 \(g\in\operatorname{SU}(4)\)，中心化子还需满足 \(\det A\,z\,w=1\)，故为
	\[
		\{(A,z,w)\in\operatorname{U}(2)\times S^1\times S^1:\det A\,z\,w=1\}.
	\]
	用 \(w=(\det A\,z)^{-1}\) 可把它参数化为 \(\operatorname{U}(2)\times S^1\)，维数为五。于是特殊酉共轭轨道维数为 \(15-5=10\)。从酉群改为特殊酉群，同时删去了作用群与稳定子中同一个一维中心方向，因而轨道维数不变。实际上两轨道作为矩阵集合也相同：任意酉共轭矩阵可乘一个适当的标量矩阵调成行列式一，而标量不会改变共轭作用。
\end{solution}

\begin{exercise}\label{g1:ex:so4-oriented-conjugacy}
	取 \(0<a<b<\pi\)，令
	\[
		g=\operatorname{diag}(R(a),R(b)),\qquad h=\operatorname{diag}(R(a),R(-b))\in\operatorname{SO}(4).
	\]
	证明 \(g,h\) 有相同的复特征值和特征多项式，并在 \(\operatorname{O}(4)\) 中共轭，但在 \(\operatorname{SO}(4)\) 中不共轭。再判断 \(g\) 与 \(g^{-1}\) 在 \(\operatorname{SO}(4)\) 中是否共轭。
\end{exercise}
\begin{solution}
	两个旋转块的复特征值分别为 \(e^{\pm\mathrm{i}a},e^{\pm\mathrm{i}b}\)，改变第二个角的符号只交换它的两特征值。因此两矩阵的特征值多重集相同。令 \(C=\operatorname{diag}(1,-1)\)，则 \(CR(b)C^{-1}=R(-b)\)，故 \(\operatorname{diag}(I_2,C)\) 给出正交共轭，其行列式为负。

	若在特殊正交群内共轭，环面内的两个元素必须在同一 Weyl 群轨道中。\(\operatorname{SO}(4)\) 的 Weyl 群只允许交换两角或同时改变两角的符号。由于 \(a,b\) 严格位于 \((0,\pi)\) 且不同，模 \(2\pi\) 后，这些操作都不能把 \((a,b)\) 变成 \((a,-b)\)。因此不存在特殊正交共轭。另一方面，\(g^{-1}\) 对应 \((-a,-b)\)，同时在两个平面上反转定向由 \(\operatorname{diag}(C,C)\in\operatorname{SO}(4)\) 实现，所以 \(g,g^{-1}\) 在特殊正交群内共轭。特征多项式没有记录实特征平面组合的这项定向信息。
\end{solution}

\begin{exercise}\label{g1:ex:sp2-centralizers}
	在 \(\operatorname{Sp}(2)\) 中，分别求下列元素的中心化子及共轭轨道实维数：
	\[
		g_1=\operatorname{diag}(e^{\mathbf i a},e^{\mathbf i b}),\quad
		g_2=e^{\mathbf i a}I_2,\quad g_3=-I_2,
	\]
	其中 \(0<a<b<\pi\)。特别说明 \(g_2\) 虽然写成标量矩阵，为何不在群的中心。
\end{exercise}
\begin{solution}
	在复矩阵模型中，\(g_1\) 的四个特征值 \(e^{\pm\mathrm{i}a},e^{\pm\mathrm{i}b}\) 两两不同。与它交换的酉矩阵为对角矩阵，再要求具有紧辛群的共轭配对形式，便只剩 \(\operatorname{diag}(e^{\mathbf i t_1},e^{\mathbf i t_2})\)。所以中心化子为二维极大环面，轨道维数为 \(10-2=8\)。

	对于 \(g_2\)，交换条件逐个四元数分量写成 \(q e^{\mathbf i a}=e^{\mathbf i a}q\)。因 \(\sin a\ne0\)，这等价于 \(q\mathbf i=\mathbf i q\)，故所有矩阵分量必须落在 \(\mathbb C=\mathbb R+\mathbb R\mathbf i\) 中。余下的方程 \(A^*A=I\) 正是 \(\operatorname{U}(2)\)，因此中心化子维数四，轨道维数六。含 \(\mathbf j\) 的矩阵一般不与 \(g_2\) 交换，所以这里的四元数标量矩阵不是中心元；只有实标量才对所有四元数分量交换。\(g_3=-I_2\) 属于中心，中心化子为整个十维群，轨道是单点。
\end{solution}

\begin{exercise}\label{g1:ex:unitary-central-cover}
	设 \(n\ge2\)，定义 \(F:S^1\times\operatorname{SU}(n)\to\operatorname{U}(n)\) 为 \(F(z,A)=zA\)。证明这是一个 \(n\) 重覆盖同态，确定其核，并说明在由圆周第一因子及行列式确定的基本群生成元下，\(F_*:\mathbb Z\to\mathbb Z\) 是乘以 \(n\)。与正文中 \(S^1\times\operatorname{SU}(n)\cong\operatorname{U}(n)\) 的微分同胚比较，解释两个结论为何没有矛盾。
\end{exercise}
\begin{solution}
	标量矩阵与所有矩阵交换，所以 \(F\) 是群同态。给定 \(U\)，选取 \(z\in S^1\) 满足 \(z^n=\det U\)，则 \(A=z^{-1}U\in\operatorname{SU}(n)\)，因此满射。核为
	\[
		\ker F=\{(z,z^{-1}I):z^n=1\}\cong\mu_n.
	\]
	单位元微分为 \((\mathrm{i}t,X)\mapsto\mathrm{i}tI+X\)。取迹可知其核为零，且两边维数都为 \(n^2\)，所以它是局部微分同胚。由\cref{g1:prop:lie-homomorphism-local-isomorphism}，\(F\) 为覆盖，纤维含 \(n\) 个点。

	源群的基本群由 \(s\mapsto(e^{2\pi\mathrm{i}s},I)\) 生成。它在 \(F\) 下变为 \(e^{2\pi\mathrm{i}s}I\)，行列式绕圆周 \(n\) 次。正文已证明行列式在基本群上为同构，故 \(F_*\) 是乘以 \(n\)。正文的微分同胚则把 \((z,A)\) 送到 \(\operatorname{diag}(z,1,\ldots,1)A\)，它是另一张映射，通常不保持直积群乘法；其圆周生成环路的行列式只绕一次。微分同胚描述底流形，中心覆盖同态描述群乘法及离散核，所保留的结构不同。
\end{solution}
